\section{The Phenomenology of the Anger Spectrum in the
Qurʾān}\label{the-phenomenology-of-the-anger-spectrum-in-the-qurux2beux101n}

\subsection{A Semantic-Network Analysis of Displeasure, Inflammation,
and Destructiveness Along an Action-Intensity
Continuum}\label{a-semantic-network-analysis-of-displeasure-inflammation-and-destructiveness-along-an-action-intensity-continuum}

\textbf{Karim Jabbary} \emph{(corresponding author)} Department of
Arabic Language and Literature, Faculty of Literature and Humanities,
Urmia University, Urmia, Iran. \emph{email:} k.jabbary@urmia.ac.ir.

\textbf{Ali Jabbary} Faculty of Computer Engineering, Urmia University,
Urmia, Iran. \emph{email:} a.jabbary@urmia.ac.ir.

(ORCID identifiers will be provided at submission.)

\begin{center}\rule{0.5\linewidth}{0.5pt}\end{center}

\subsection{Abstract}\label{abstract}

The Qurʾān deploys an exceptionally articulated lexicon for the
\emph{anger spectrum}, yet existing scholarship has treated it either
word-by-word or through normative-ethical frames, leaving the
structural-relational dimension of the field unexplored. This study
advances the thesis that fourteen core Arabic roots --- \emph{ʾff, krh,
ḍyq, ḥzn, ʾsf, nqm, sxṭ, mqt, ġḍb, ḥrd, ġyẓ, bġy, ṭġy, ʿtw} --- form a
graded semantic field organised by \emph{intensity of action} across six
phenomenological stages: pre-anger displeasure, inner pressure and
contraction, evaluative aversion, active anger, compressed/explosive
rage, and behavioural outcomes (with caveats). We integrate three
methodological lenses --- Izutsu's semantic-field approach,
Kennedy-Sassoon scalar semantics, and Lakoff-Kövecses conceptual
metaphor theory --- and validate the continuum against the Quranic
Arabic Corpus (Dukes 2011) through a fully reproducible Python pipeline
analysing 312 attestations across the 6,236-verse corpus. We report
three principal findings under uncertainty-aware quantitative reporting
(effect sizes, FDR control, permutation null, bootstrap CIs). First, the
asymptotic χ²(5) = 227.15 (Cohen's \emph{w} = 0.85, large) is large, but
a marginal-preserving permutation null yields \emph{p} = 0.24, so the
omnibus is treated as \emph{consistent with} but not confirmatory of
non-uniform stage distribution; the substantive case rests on
qualitative lexicographic, exegetical, and metaphorical evidence. Three
monotonic-trend tests fail to reject (Spearman ρ = −0.09; Mann-Kendall
\emph{p} = 0.71; CA-OLS \emph{p} = 0.40) --- the expected result, since
the continuum is an ordering of semantic intensity, not frequency.
Second, \emph{sakhaṭ} attests exclusively in Medinan suras (0/4; raw
\emph{p} = 0.005, robust across baselines 0.70--0.78; Holm-Bonferroni
adjusted \emph{p} = 0.05 --- borderline post-FDR significance), proposed
as the only confirmed instance of a more general phenomenon of
\emph{discursive-contextual lexical specialisation}; \emph{asaf} (5/0
Meccan), \emph{naqm} (12/5 Meccan), and \emph{ʿutuww} (9/1 Meccan) are
reported as contextual observations not surviving FDR correction. Third,
network-centrality analysis identifies \emph{baghy} as a bridging node
(point-estimate betweenness 0.15, closeness 0.55) between the anger
spectrum and the broader Qurʾānic moral-evaluation field, with bootstrap
95\% CIs reported transparently. We argue that Q. 67:8 (\emph{takādu
tamayyazu min al-ġayẓ}) furnishes an \emph{unusually transparent}
instance of ``container collapse'' --- the verb \emph{tamayyaza}
lexicalising the disintegration of the vessel itself rather than the
exit of pressurised contents --- softened by the recognition that no
comparative cross-cultural survey of historical anger-vessel imagery has
yet been conducted. A κ = 0.79-validated tafsir-coding of the 145
Stage-6 attestations supports a \emph{bidirectional-causation} reading:
\emph{baghy}, \emph{ṭughyān}, and \emph{ʿutuww} are not exclusively
anger-derived but stand in mutual causal relation to the anger field.
The findings have implications for Qurʾānic translation (conditional on
follow-up evaluation) and for culturally-grounded Islamic-psychology
interventions.

\textbf{Keywords:} Qurʾān, anger spectrum, semantic field, cognitive
linguistics, scalar semantics, conceptual metaphor, emotion regulation,
Arabic corpus linguistics, Izutsu, Islamic psychology.

\begin{center}\rule{0.5\linewidth}{0.5pt}\end{center}

\subsection{1. Introduction}\label{introduction}

\subsubsection{1.1. Statement of the
problem}\label{statement-of-the-problem}

The Qurʾānic lexicon for the \emph{anger spectrum} exhibits a precision
that survives, only with difficulty, the attempts of translation to
render it into modern languages. Where English versions level
\emph{uff}, \emph{karh}, \emph{ḍīq}, \emph{ḥuzn}, \emph{asaf},
\emph{naqm}, \emph{sakhaṭ}, \emph{maqt}, \emph{ghaḍab}, \emph{ḥard},
\emph{ghayẓ}, \emph{baghy}, \emph{ṭughyān}, and \emph{ʿutuww} into a
small set of broad equivalents---chiefly ``displeasure,'' ``distress,''
``grief,'' ``anger,'' ``rage,'' ``transgression,'' and
``rebellion''---the Arabic original distinguishes them with the
precision of a graded taxonomy of irritation, aversion, anger, and its
outcomes. The hypothesis advanced in this paper is structural: that
fourteen of the Qurʾān's thematically central anger-spectrum roots are
not synonyms or scattered semantic singletons but coordinated nodes in a
single \emph{graded} semantic field, organised along the dimension of
\emph{intensity of action} (\emph{shiddat al-kunish} in the Persian
formulation of the original idea). The field opens in pre-anger
displeasure (the murmured \emph{uff}, the inner \emph{karāha}), passes
through inner pressure and contraction, evaluative aversion, active
anger, and compressed/explosive rage, and---at the behavioural
pole---issues into outcomes (\emph{baghy}, \emph{ṭughyān},
\emph{ʿutuww}) whose causation, as we shall argue, is bidirectional
rather than purely emotional. We treat the fourteen roots as discrete
attestation points on this continuum and validate the proposed ordering
through systematic corpus-based analysis.

This claim, if it holds, has consequences far beyond philology. It
implies that the Qurʾān encodes---through the architecture of its
lexicon, not through any explicit theoretical statement---a
sophisticated implicit theory of the genealogy of moral failure: a
theory according to which the social-structural pathologies the text
most frequently denounces (transgression, tyranny, defiance) are
\emph{related to}, and in many cases developmentally downstream from,
inner experiences of irritation, dislike, constriction, and grief. As we
shall argue in §4 and §5, the relationship at the behavioural pole is
bidirectional: anger may produce
\emph{baghy}/\emph{ṭughyān}/\emph{ʿutuww}, but these behaviours can
equally produce \emph{anger} (in their victims, in the prophets, in the
divine response). The distance between the murmured \emph{uff} of Q.
17:23 and the \emph{ʿutuww} of Q. 25:21 is not categorical but graded;
it is the distance walked, in the absence of intervention, by an
unmanaged life of irritation, dislike, and anger.

\subsubsection{1.2. Aims and originality}\label{aims-and-originality}

This paper has three aims. The first is \textbf{methodological}: to
demonstrate that combining Izutsu's qualitative semantic-field approach,
scalar-semantic ordering in the Kennedy--Sassoon tradition, and
conceptual-metaphor theory in the Lakoff--Kövecses tradition produces
stronger philological claims than any of the three frameworks in
isolation. To our knowledge---and we say this advisedly, given that the
Persian-language Qurʾānic-linguistic literature in particular is not
exhaustively indexed in international databases---this trifold
integration is here brought together in Qurʾānic studies for the first
time as a unified analytic framework over a corpus pipeline. We
acknowledge that Qāʾimīniyā (1390/2011, 1393/2014) and Pakatchi already
integrate Izutsu's field method with cognitive-linguistic
conceptual-metaphor analysis, and that Persian-language scalar-semantic
work on Qurʾānic vocabulary may exist in venues we have not been able to
canvass; the \emph{novelty} of the present paper is therefore best
stated as the \emph{systematic combination} of the three frameworks
within a single computational-empirical pipeline, not as the
introduction of any one framework. The second aim is
\textbf{descriptive}: to map every Qurʾānic occurrence of the fourteen
core roots, to render that map publicly verifiable through an open
concordance, and to extract from it empirical regularities that previous
scholarship---working in either tafsir, lexicography, or single-word
semantic analysis---was not in a position to detect. The third aim is
\textbf{interpretive}: to argue that the resulting continuum constitutes
a coherent ``logic of anger escalation'' with deep structural affinities
to modern psychology of emotion, and to identify the specific points at
which the Qurʾānic continuum \emph{exceeds} the explanatory range of the
contemporary models.

The paper's claim to \textbf{originality} is best stated by
disaggregating three logically distinct contributions, with ``to our
knowledge'' qualifications throughout. \textbf{(i) Methodological
synthesis} (moderate): the trifold integration of Izutsu's field method,
Kennedy--Sassoon scalar semantics, and Lakoff--Kövecses CMT under a
corpus-empirical pipeline. We acknowledge Qāʾimīniyā (1390/2011,
1393/2014) and Pakatchi-Afrashi (1399/2020) already integrate Izutsu and
CMT in Persian-language Qurʾānic scholarship; the specifically novel
element is the Kennedy--Sassoon scalar-semantic layer within the
pipeline. \textbf{(ii) Empirical findings} (variable strength, reported
transparently): of four distributional patterns observed, only the
exclusively-Medinan \emph{sakhaṭ} (0/4; raw \emph{p} = 0.005, robust
across baselines 0.70--0.78; \textbf{Holm-Bonferroni adjusted \emph{p} =
0.05}) survives FDR control at α = 0.05 --- borderline rather than
emphatic. The \emph{asaf} (5/0), \emph{naqm} (12/5), and \emph{ʿutuww}
(9/1) tilts do not survive correction and are reported as contextual
observations. That one of fourteen pre-specified tests survives
Holm-Bonferroni is the kind of result honest corpus work should report.
\textbf{(iii) Interpretive contribution} (the most substantive): the
\emph{bidirectional-causation reading of Stage 6}, operationalised
through a κ = 0.79-validated tafsir-coding (anger-derived \emph{A} =
26\%; structurally-derived \emph{S} = 49\%; joint \emph{A+S} = 25\%;
full coding at \texttt{data/concordance/stage6\_causation\_coding.csv}).
Re-running the χ²(1) phenomenology-vs-outcomes split on the \emph{A}
subset alone yields χ²(1) = 81.4 (\emph{p} \textless{} 0.001), with the
phenomenological core dominating anger-derived outcomes
\textasciitilde4.4×. This reconciles a puzzling distributional finding
with the structurally-driven character of \emph{baghy}, \emph{ṭughyān},
\emph{ʿutuww} in classical exegesis. To these we add \textbf{(iv)} a
fully reproducible MIT-licensed pipeline (SHA-256 pinning, deterministic
permutation/bootstrap, 50/50 manual validation) and \textbf{(v)} the
comparative observation that Q. 67:8 (\emph{takādu tamayyazu min
al-ġayẓ}) is an unusually transparent instance of ``container
collapse'', softened in §5.6 by the absence of a cross-cultural survey
of historical anger-vessel imagery.

\subsubsection{1.3. Prior work and the
gap}\label{prior-work-and-the-gap}

Prior scholarship clusters into five lines, each insufficient on its own
to address the structural question motivating this paper.
\textbf{Single-word lexical studies}: the most directly relevant
Persian-language precedent --- Narimani, Eghbali, and Chari (2021;
\emph{JQST} 54:1) on \emph{ghayẓ}, \emph{ghaḍab}, and \emph{sakhaṭ} ---
establishes on tafsir grounds that \emph{ghaḍab} is central,
\emph{sakhaṭ} reserved for divine displeasure with hypocrites, and
\emph{ghayẓ} exclusive to unbelievers/hypocrites in human-attribution
contexts; we confirm and extend their findings with
corpus-distributional evidence and the \emph{kaẓm} analysis. Qarizadeh
et al.~(2024) on \emph{ṭughyān} shows the root attests in 27 suras with
the central meaning of transgressing the limit. Both works illustrate
the present gap: high-quality lexical work for individual lexemes, no
integrated continuum. \textbf{Ethics-of-anger studies} focus on
\emph{kaẓm al-ghayẓ} (Q. 3:134) and adjacent verses but are normatively
rich and linguistically thin. \textbf{Cognitive-semantic studies} in the
Izutsu tradition --- Qāʾimīniyā (1390/2011, 1393/2014); Pakatchi-Afrashi
(1399/2020); Albayrak (2020) --- have applied cognitive linguistics to
Qurʾānic vocabulary but not treated the anger spectrum as a graded
continuum. The cross-linguistic CMT literature on anger is decisively
updated by \emph{Metaphors of ANGER across Languages} (Kövecses, Benczes
\& Szelid eds., 2025; De Gruyter, 2 vols.); Classical Arabic is
conspicuously not separately treated in Vols. 1--2, a gap this paper
begins to fill. \textbf{Translation-criticism studies} are quantitative
at scale in 2025: Gaanoun \& Alsuhaibani (2025; \emph{Nature HSSC}) show
inconsistent sentiment polarity across seven English translations;
Alsuhaibani et al.'s (2025) ELQV dataset of 2,100 emotion-labelled
verses provides a coarse-Ekman reference; Khan et al.~(2025;
\emph{Frontiers Signal Proc.}) survey the field; Cogent A\&H 2025
catalogues flattening risks. None of these supplies a \emph{graded
reference model of the lexical spectrum} at the lexeme level.
\textbf{Islamic psychology} (e.g.~\emph{IJCP} 2023 on
fear/anger/sadness/shame in Qurʾān-and-ḥadīth) tends to draw on
\emph{aḥādīth} and leave systematic linguistic analysis of the Qurʾānic
surface in the background. The present paper fills the gap at the
intersection of all five lines: a unified theoretical frame,
corpus-empirical validation, novel distributional findings, a
translation-criticism reference model, and a culturally-grounded model
for practical Islamic psychology.

\subsubsection{1.4. Research questions and
hypotheses}\label{research-questions-and-hypotheses}

We address five research questions: (RQ1) the internal structure of the
Qurʾānic anger-spectrum field and its principal nodes; (RQ2) the
continuum ordering by action-intensity and the theoretical framework
legitimating it; (RQ3) corpus-distributional corroboration of the
continuum (frequency, Meccan/Medinan, morphology, network centrality);
(RQ4) structural homology with Gross, Plutchik, Spielberger, and
Lazarus-Scherer appraisal theory; (RQ5) practical implications for
Qurʾānic translation and clinical Islamic psychology.

The principal \textbf{Qurʾānic Action-Intensity Continuum Hypothesis}
holds that the fourteen core roots (ʾff, krh, ḍyq, ḥzn, ʾsf, nqm, sxṭ,
mqt, ġḍb, ḥrd, ġyẓ, bġy, ṭġy, ʿtw) operate as fourteen distinct grades
on a single semantic continuum structured by \emph{intensity of action},
organised in six phenomenological stages, supported by qualitative
(etymological, exegetical, metaphorical) and quantitative (frequency,
distribution, network) evidence. Five subsidiary hypotheses: (H1)
non-uniform six-stage frequency distribution; (H2) near-equal split
between phenomenological core (S1--5) and behavioural pole (S6),
supporting the bidirectional Stage-6 reading; (H3)
discursively-specialised exclusive-Medinan distribution of
\emph{sakhaṭ}; (H4) \emph{baghy} as bridge node in the co-occurrence
network; (H5) Q. 67:8 as transparent instance of ``container collapse''
CMT failure mode.

\begin{center}\rule{0.5\linewidth}{0.5pt}\end{center}

\subsection{2. Theoretical Framework}\label{theoretical-framework}

\subsubsection{2.1. Izutsu's semantic-field
method}\label{izutsus-semantic-field-method}

Toshihiko Izutsu, in three foundational works---\emph{The Structure of
the Ethical Terms in the Koran} (1959), \emph{God and Man in the Koran}
(1964), and \emph{Ethico-Religious Concepts in the Qurʾān}
(1966/2002)---introduced a methodological revolution to Qurʾānic
studies. Drawing on the structuralist linguistics of the
Trier--Weisgerber tradition and on East Asian philosophy of language,
Izutsu argued that the central terms of the Qurʾān---\emph{Allāh},
\emph{īmān}, \emph{kufr}, \emph{ẓulm}, \emph{taqwā}---are not isolated
propositions but nodes in a network of syntagmatic, paradigmatic, and
oppositional relations. The aggregate of these relations constitutes
what he called the ``Qurʾānic \emph{Weltanschauung}.''

Izutsu's analytic vocabulary distinguishes three layers: \emph{focus
words} are the central organising concepts of a given semantic field;
\emph{semantic fields} are the networks of mutually defining lexemes
formed around focus words; and \emph{key words} are the terms that recur
across multiple fields and serve to integrate them into a single
worldview. This three-layer structure has proved extraordinarily
productive in subsequent Qurʾānic semantic work, both in English (Saleh,
Rippin) and Persian (Qāʾimīniyā 2011, 2014; Pakatchi).

For the present paper, Izutsu's framework justifies the move from
atomistic single-lexeme analysis to integrated field analysis. The
fourteen target roots are not chosen at random but as the field around
the focus words \emph{ghaḍab} (the central anger term) and \emph{baghy}
(the central behavioural-outcome term), with the lower-intensity cluster
(\emph{uff}, \emph{karh}, \emph{ḍīq}, \emph{ḥuzn}, \emph{asaf})
identifying the pre-anger and inner-pressure boundary of the field and
\emph{ʿutuww} its upper-intensity boundary. Yet Izutsu's framework, in
its original form, is silent on a question central to our enterprise:
granted that these lexemes co-belong to a field, \emph{what is the
structural relation between them within the field?} It is at this point
that we require scalar semantics.

\subsubsection{2.2. Scalar semantics: from network to
continuum}\label{scalar-semantics-from-network-to-continuum}

The contemporary linguistic theory of \emph{gradable predicates}
(Kennedy 1999, 2007; Sassoon 2010) supplies the missing apparatus. A
graded predicate maps to a scale characterised by three components: a
\emph{dimension} of measurement, an \emph{ordering} on that dimension,
and a unit of comparison. Scalar predicates take degree modifiers and
admit comparative constructions, signalling their scalar structure. The
Qurʾānic anger field is precisely such a scale: the dimension is
\emph{intensity of action} (decomposable into locus, experiential
intensity, scope of consequence); the ordering follows the six-stage
progression; the implicit unit is the \emph{degree of externalisation}
of emotion from inner phenomenology to social rupture. Without this
framework, the claim that \emph{ghayẓ} is ``more intense than''
\emph{ghaḍab} relies on tacit native-speaker intuition.

\subsubsection{2.3. Conceptual metaphor theory and the container
schema}\label{conceptual-metaphor-theory-and-the-container-schema}

Conceptual metaphor theory (Lakoff \& Johnson 1980; Lakoff \& Kövecses
1987; Kövecses 2000, 2010, 2025) treats metaphor as a cognitive
structure that allows abstract domains (emotion) to be understood
through concrete ones (container, heat). Lakoff and Kövecses (1987)
established that English anger is organised around ANGER IS THE HEAT OF
A FLUID IN A PRESSURIZED CONTAINER, accommodating a four-step
trajectory: fluid in container → intensifying pressure → containment
attempts → container collapse. Our central observation is that the
Qurʾānic triplet \emph{ghayẓ--kaẓm--tamayyuz} operationalises this exact
structure thirteen centuries early. \emph{Ghayẓ} is glossed by all four
authoritative lexica as ``compressed contained anger'' --- fluid in
sealed vessel. \emph{Kaẓm}, etymologically tying off a waterskin, is the
agent-internal regulation that maintains containment (Q. 3:134
\emph{al-kāẓimīn al-ghayẓ}). \emph{Tamayyuz} (Q. 67:8) depicts the
failure mode: the container itself bursts apart. The Qurʾānic verb
derives from \emph{m-y-z} (to part, sever) and figures the
disintegration of the \emph{vessel}, not (as in English \emph{explode},
\emph{burst}, \emph{blow up}) the \emph{exit} of contents --- an
unusually transparent instantiation of the failure-mode schema. We treat
\emph{tamayyuz} as a \emph{manifestation} of Stage 5, not an independent
core root.

\subsubsection{2.4. The integrated
framework}\label{the-integrated-framework}

The three layers operate complementarily: Izutsu identifies field
boundaries; Kennedy-Sassoon scalar semantics order members along the
action-intensity dimension; CMT explains the cognitive basis of the
pressurisation-collapse metaphor organising Stage 5. Each supplies what
the others lack --- without Izutsu, no motivation for
field-coordination; without scalar semantics, no apparatus for ordering;
without CMT, no cognitive grounding for the climactic
\emph{kaẓm--tamayyuz} sequence.

\begin{center}\rule{0.5\linewidth}{0.5pt}\end{center}

\begin{figure}
\centering
\includegraphics[width=1\linewidth,height=\textheight,keepaspectratio]{fig1_continuum.pdf}
\caption{The six-stage Qurʾānic anger-spectrum continuum, mapped along a
single horizontal axis of \emph{action-intensity}. Each node corresponds
to one of the fourteen core Arabic roots; node area is scaled by
attestation frequency in the Quranic Arabic Corpus (Dukes 2011). Stages
cluster the lexemes by their dominant phenomenological signature: Stage
1 (pre-anger displeasure, \emph{ʾff--krh}) marks vocal/inner
displeasure; Stage 2 (inner pressure \& contraction,
\emph{ḍyq--ḥzn--ʾsf}) marks pre-anger affective contraction and grief;
Stage 3 (evaluative aversion, \emph{nqm--sxṭ--mqt}) marks
moral-judgmental rejection; Stage 4 (active anger, \emph{ġḍb--ḥrd})
marks outward-oriented anger; Stage 5 (compressed/explosive rage,
\emph{ġyẓ} with \emph{tamayyuz} as manifestation) marks containment and
its failure; Stage 6 (behavioural outcomes with caveats,
\emph{bġy--ṭġy--ʿtw}) marks the externalised structural outcome whose
causal relation to anger is, as we argue in §4.6 and §5,
\emph{bidirectional} rather than purely downstream. The visual mass at
Stage 6 (n = 145) is comparable to the combined mass of Stages 1--5 (n =
167); the resulting near-balance (χ²(1) = 1.55, fails to reject) is
theoretically consistent with the bidirectional reading.}
\end{figure}

\subsection{3. Methods}\label{methods}

\subsubsection{3.1. Corpus}\label{corpus}

The empirical analysis uses the \textbf{Quranic Arabic Corpus} (QAC),
version 0.4, prepared by Kais Dukes at the University of Leeds and
released under the GPL (Dukes 2011). The QAC provides a morphologically
tagged record for each of the 128,219 segmentable units in the Qurʾānic
text. Every record specifies: (a) location in the form
\emph{(sūra:āya:word:segment)}; (b) surface form in Buckwalter ASCII
transliteration; (c) part-of-speech tag; (d) morphological features
including root, lemma, gender, number, case, and (where applicable)
semantic flags. Verse text in Uthmanic \emph{rasm} and Meccan/Medinan
classification of suras are drawn from the Tanzil project (Tanzil 2008),
distributed under CC-BY-ND 3.0.

\subsubsection{3.2. Selection of target
roots}\label{selection-of-target-roots}

The fourteen core roots were selected by three concurrent criteria,
applied transparently and reported here so that the inventory is not
vulnerable to the charge of post-hoc selection: (i) \emph{exegetical
centrality} --- appearance as a focal reference in at least three of the
four authoritative lexica (Ibn Fāris's \emph{Maqāyīs}, al-Rāghib's
\emph{Mufradāt}, Ibn Manẓūr's \emph{Lisān}, and Mostafavi's
twentieth-century \emph{Taḥqīq}); (ii) \emph{minimum corpus frequency}
of one attestation, admitting Qurʾānic hapax legomena (\emph{ḥard}, Q.
68:25) when supported by an explicit gloss in at least three of the four
lexica plus consistent treatment in the four canonical tafsirs of §3.5;
(iii) \emph{thematic membership} in the anger-spectrum field without
significant overlap with adjacent fields (epistemology, faith, normative
ethics).

\textbf{Excluded candidate roots} (transparency): \emph{q-h-r}
(dominance/power, \emph{al-Qahhār} divine name); \emph{ḥ-n-q}
(post-Qurʾānic, ḥadīth-only); \emph{ʿ-n-f} (adverbial harshness, not
emotion); \emph{ḍ-r-b} (polysemous, conflict-scene);
\emph{q-ṣ-w}/\emph{q-s-y} (hardness-of-heart, epistemic-moral);
\emph{gh-l-ẓ} (adverbial harshness); \emph{ʾ-n-f} (nominal \emph{anf}
``nose'', post-Qurʾānic anger sense); \emph{š-n-ʾ} (binary
loving/hating, reclassified as umbrella). \emph{Karh} (Stage 1) is
borderline --- \emph{karāha} is affective, \emph{ikrāh} (coercion)
structural; we include the affective sense and discuss the structural
extension separately (§4.3.1). Five umbrella moral terms (\textbf{šnʾ,
jrm, fsq, ẓlm, ʿdw}) are used for robustness (§4.9) and cross-field
co-occurrence (Table 4), excluded from the principal continuum either
because of inflated cardinality (\emph{ẓlm} n=315 functions as a
meta-category) or because they belong to adjacent normative-ethics
fields.

\subsubsection{3.3. Computational
pipeline}\label{computational-pipeline}

The Python (MIT-licensed) pipeline operates in five steps: parse the QAC
morphology file; filter on target roots; collapse multi-segment
morphemes to surface words preserving the \emph{(sūra:āya)} coordinate;
join with sura-level metadata for Meccan/Medinan classification; emit
master concordance, per-root concordances, distributional summaries, and
network/centrality CSVs. The complete analysis is reproducible from a
fresh clone of the repository in four commands (\texttt{pip\ install},
\texttt{extract\_concordance.py}, \texttt{advanced\_metrics.py},
\texttt{visualize\_*.py}).

\textbf{Reproducibility audit (summary).} QAC v0.4 (Dukes, May 2011) is
pinned by SHA-256 in the repository; minimum-version dependencies are
listed in \texttt{requirements.txt}; outputs are released under MIT
(code) and CC-BY 4.0 (manuscripts/figures); the 14 headline exemplar
verses and a stratified 50-segment audit on the four most frequent core
roots (\emph{ḥzn}, \emph{bġy}, \emph{ġḍb}, \emph{ṭġy}) returned 100\%
correct segment-to-root assignment against Tanzil. Full audit at
\texttt{data/concordance/validation\_audit.md}.

\subsubsection{3.4. Quantitative analyses}\label{quantitative-analyses}

Quantitative tests beyond frequency tabulation: (a) χ²(5) and χ²(1) for
uniform-stage and S1--5/S6 nulls, with Cramér's \emph{V} and Cohen's
\emph{w}, plus a marginal-preserving permutation null (10,000 perms,
deterministic seed) for the omnibus; (b) three monotonic-trend tests
(Spearman ρ, Mann-Kendall, Cochran-Armitage-style OLS slope); (c)
two-sided exact binomials with the corpus-derived 0.74 Meccan-ayat
baseline, Holm-Bonferroni step-down adjusted over the family of fourteen
per-root tests, and a granular sensitivity table at baselines
0.70--0.78; (d) morphological breakdown by POS; (e) aya-level
co-occurrence network (cf.~Mottaghi et al.~2024) with weighted degree,
betweenness, closeness, and eigenvector centrality, accompanied by
non-parametric bootstrap 95\% percentile CIs (1,000 ayat-resamples with
replacement). The combined uncertainty-aware reporting is, to our
knowledge, the most rigorous yet applied to Qurʾānic
computational-corpus studies; the wider 2025 cross-scripture NLP
literature (e.g.~Nandan et al.~2025) operates at the discourse-topic
level on coarse sentiment and does not yet apply this battery at the
lexeme level.

\subsubsection{3.5. Qualitative analyses}\label{qualitative-analyses}

Qualitative complements: (a) comparative etymology across the four
authoritative lexica (Maqāyīs, Mufradāt, Lisān, Taḥqīq); (b) comparative
\emph{tafsīr} analysis on a broadened set of nine commentaries --- the
four primary (al-Ṭabāṭabāʾī's \emph{al-Mīzān}, al-Zamakhsharī's
\emph{al-Kashshāf}, al-Ṭabarsī's \emph{Majmaʿ al-Bayān}, Ibn ʿĀshūr's
\emph{al-Taḥrīr}) supplemented with al-Ṭabarī, al-Qurṭubī, Ibn Kathīr,
Sayyid Quṭb, and ʿAbduh-Riḍā's \emph{al-Manār}; for Sufi-tradition
cross-checking on \emph{kaẓm al-ghayẓ} (§4.7), al-Qushayrī's
\emph{Laṭāʾif al-Ishārāt}. Where the broader set diverges from the
primary four on key claims (whether \emph{ghayẓ} is exclusively
predicated of unbelievers; whether \emph{sakhaṭ} is exclusively directed
at the \emph{Munāfiqūn}; whether \emph{naqm} admits a non-retributive
sense), divergences are recorded in §4 and corresponding claims
qualified. (c) phenomenological analysis of contextual markers
(\emph{ḍīq al-ṣadr}, \emph{ḥuzn al-qalb}, \emph{asaf al-nafs},
\emph{ḍāqa dharʿan}); (d) illustrative translation-criticism on five
English translations (Yusuf Ali, Pickthall, Arberry, Saheeh
International, Hilali-Khan) at canonical exemplar verses (§4.10) ---
framed throughout as illustrative rather than systematic.

\subsubsection{3.6. Validity check}\label{validity-check}

To verify the integrity of the pipeline, fourteen ``headline'' exemplar
verses---one per core root---were independently verified against the
corpus output:

\begin{itemize}
\tightlist
\item
  Q. 17:23 (\emph{uff}); Q. 2:216 (\emph{karh}); Q. 15:97 (\emph{ḍyq});
  Q. 2:38 (\emph{ḥzn}); Q. 43:55 (\emph{asaf}); Q. 7:126 (\emph{naqm});
  Q. 47:28 (\emph{sakhaṭ}); Q. 40:10 (\emph{maqt}); Q. 48:6
  (\emph{ghaḍab}); Q. 68:25 (\emph{ḥard}); Q. 3:134 (\emph{ghayẓ}; with
  manifestation Q. 67:8 \emph{tamayyuz}); Q. 42:42 (\emph{baghy}); Q.
  96:6 (\emph{ṭughyān}); Q. 25:21 (\emph{ʿutuww}).
\end{itemize}

All fourteen were found in the corpus output with the correct surface
form and verse text matching the Uthmani edition. This validation
confirms that the Buckwalter transliteration and segment-collapse
algorithms operate without systematic error.

\begin{center}\rule{0.5\linewidth}{0.5pt}\end{center}

\subsection{4. Findings}\label{findings}

\subsubsection{4.1. The four-lexicon etymology
table}\label{the-four-lexicon-etymology-table}

Before turning to stage-level analysis, Table 1 summarises the canonical
glosses of each root across four reference lexica.

\textbf{Table 1. Comparative etymology of the fourteen core roots across
the four authoritative Arabic lexica}

\begin{longtable}[]{@{}
  >{\centering\arraybackslash}p{(\linewidth - 10\tabcolsep) * \real{0.1200}}
  >{\centering\arraybackslash}p{(\linewidth - 10\tabcolsep) * \real{0.2400}}
  >{\raggedright\arraybackslash}p{(\linewidth - 10\tabcolsep) * \real{0.1600}}
  >{\raggedright\arraybackslash}p{(\linewidth - 10\tabcolsep) * \real{0.1600}}
  >{\raggedright\arraybackslash}p{(\linewidth - 10\tabcolsep) * \real{0.1600}}
  >{\raggedright\arraybackslash}p{(\linewidth - 10\tabcolsep) * \real{0.1600}}@{}}
\toprule\noalign{}
\begin{minipage}[b]{\linewidth}\centering
Stage
\end{minipage} & \begin{minipage}[b]{\linewidth}\centering
Root
\end{minipage} & \begin{minipage}[b]{\linewidth}\raggedright
Maqāyīs (Ibn Fāris)
\end{minipage} & \begin{minipage}[b]{\linewidth}\raggedright
Mufradāt (al-Rāghib)
\end{minipage} & \begin{minipage}[b]{\linewidth}\raggedright
Lisān al-ʿArab (Ibn Manẓūr)
\end{minipage} & \begin{minipage}[b]{\linewidth}\raggedright
Taḥqīq (Mostafavi)
\end{minipage} \\
\midrule\noalign{}
\endhead
\bottomrule\noalign{}
\endlastfoot
1 & ʾ-f-f & ``kalimatu taḍajjurin'' (a word of vexation); the wretched
of fingernail-detritus etymologically; smallest audible expression of
displeasure & the lowest verbal expression of displeasure (\emph{adnā
kalimati al-taḍajjur}) & term of irritation (\emph{taḍajjur}) directed
at someone or something disliked; cf.~\emph{uffah} (sub-vocal sigh) &
minimal vocal sign of inner aversion; pre-anger affective register \\
1 & k-r-h & dislike, finding burdensome; coercion (\emph{ikrāh}) &
aversion of the soul; structurally extended to coercion & what is
disliked by nature or by reason & inner \emph{karāha} vs.~structural
\emph{ikrāh} (the latter is not an emotion) \\
2 & ḍ-y-q & constriction; absence of opening & the state of being
pressed upon, contrasted with capacity (\emph{saʿah}) & reduction of
available space, from \emph{saʿah} to \emph{ḍīq} & inner or outer
absence of expansion (\emph{basṭ}) \\
2 & ḥ-z-n & rough/uneven ground → heaviness of soul & sorrow as the
opposite of joy (\emph{faraḥ}) & rugged ground, then roughness of the
self & affective response to perceived loss \\
2 & ʾ-s-f & intense grief mixed with anger & extreme anger conjoined
with grief & reaching the limit of grief and anger & composite of grief
and proximate anger \\
3 & n-q-m & \emph{al-naqimah}: disapproval (\emph{inkār}) coupled with
intent to punish; cf.~\emph{intaqama min} = exact retribution from &
objecting to (\emph{ankara ʿalā}) and seeking redress (\emph{ʿuqūbah})
against a reprehensible act & strong denunciation (\emph{ankara ashadda
al-inkār}) yoked to \emph{ʿuqūbah}; basis of divine attribute
\emph{al-Muntaqim} & evaluative anger oriented toward retribution;
cognition-grounded, action-tilted \\
3 & s-kh-ṭ & severe aversion and ethical rejection & the opposite of
\emph{riḍā} (consent), with intensity & severe non-acceptance with
judgment & declarative rejection grounded in moral assessment \\
3 & m-q-t & the most intense form of \emph{ibghāḍ} (Zajjāj) & aversion
arising from observing a \emph{qabīḥ} act (\emph{al-Rāghib}) &
\emph{ashaddu al-ibghāḍ}; aversion-distance rather than
approach-aggression & purely evaluative moral aversion \\
4 & gh-ḍ-b & inner agitation and motion & the boiling of the heart's
blood seeking retribution & a rising perturbation of the soul & directed
emotion enabling retributive action \\
4 & ḥ-r-d & \emph{ḥarada}: to be angry; also the act of \emph{qaṣd}
(deliberate aim); the lexicon couples affect and intent & conjoining
\emph{ghaḍab} with \emph{qaṣd}: anger directed by purpose, viz.~denial &
\emph{ḥarada} ``to act/move toward an aim with resolve while in anger'';
cf.~\emph{ḥarīd} = isolated, the angry-aloof & anger weaponised through
deliberate resource-denial (Q. 68:25 \emph{aṣḥāb al-jannah}); affective
+ volitional + miserly composite \\
5 & gh-y-ẓ & pressure and surge, contained & the most intense forms of
contained anger & filling of the self with restrained anger &
high-pressure emotional accumulation \\
6 & b-gh-y & seeking beyond rightful bounds & overstepping the path of
right & dominance through aggression & aggressive action with intent to
oppress \\
6 & ṭ-gh-y & overflow of water past its bank & crossing the permitted
measure & severe transgression of the limit & rebellious shattering of
normative bounds \\
6 & ʿ-t-w & severe rebellion with arrogance & extreme
self-aggrandisement in disobedience & obstinate excess attached to
\emph{istikbār} & entrenched rebellious posture (a property of moral
character) \\
\end{longtable}

The table shows that the ordering from \emph{uff} to \emph{ʿutuww} is
not the construction of the present analysis but is internally encoded
in the lexicographic tradition itself: the four authoritative lexica
already gloss each root with progressively more externalised and
structural meanings, ascending from the murmured \emph{uff}, through
inner \emph{karāha}, \emph{ḍīq}, and \emph{ḥuzn}, through evaluative
\emph{naqm}/\emph{sakhaṭ}/\emph{maqt}, to active
\emph{ghaḍab}/\emph{ḥard}, to compressed \emph{ghayẓ}, and into the
externalised behavioural pole at
\emph{baghy}/\emph{ṭughyān}/\emph{ʿutuww}.

\subsubsection{4.2. Distributional
summary}\label{distributional-summary}

Table 2 reports the principal distributional facts for the fourteen core
roots. The table integrates frequency, Meccan/Medinan split, and (for
the roots with sufficient morphological variation) morphological
profile; it forms the empirical backbone of the analyses that follow.

Figures 2 and 3 visualise the distributional profile graphically: Figure
2 reports per-root frequency and the per-stage stage-totals; Figure 3
plots the Meccan/Medinan distribution and highlights the four
discursively-specialised roots (\emph{sakhaṭ}, \emph{asaf}, \emph{naqm},
\emph{ʿutuww}) discussed below. Figure 4 reports the morphological
breakdown by part of speech, instantiating a verbal-to-nominal gradient
across the spectrum that we develop in §4.6.1 below.

\textbf{Table 2. Distributional profile of the fourteen core roots}

\begin{longtable}[]{@{}
  >{\centering\arraybackslash}p{(\linewidth - 10\tabcolsep) * \real{0.1667}}
  >{\centering\arraybackslash}p{(\linewidth - 10\tabcolsep) * \real{0.1667}}
  >{\centering\arraybackslash}p{(\linewidth - 10\tabcolsep) * \real{0.1667}}
  >{\centering\arraybackslash}p{(\linewidth - 10\tabcolsep) * \real{0.1667}}
  >{\centering\arraybackslash}p{(\linewidth - 10\tabcolsep) * \real{0.1667}}
  >{\raggedright\arraybackslash}p{(\linewidth - 10\tabcolsep) * \real{0.1667}}@{}}
\toprule\noalign{}
\begin{minipage}[b]{\linewidth}\centering
Stage
\end{minipage} & \begin{minipage}[b]{\linewidth}\centering
Root
\end{minipage} & \begin{minipage}[b]{\linewidth}\centering
Total
\end{minipage} & \begin{minipage}[b]{\linewidth}\centering
Meccan
\end{minipage} & \begin{minipage}[b]{\linewidth}\centering
Medinan
\end{minipage} & \begin{minipage}[b]{\linewidth}\raggedright
Notes
\end{minipage} \\
\midrule\noalign{}
\endhead
\bottomrule\noalign{}
\endlastfoot
1 & ʾff (أف) & 3 & 3 & 0 & All Meccan; lowest tier of verbal
displeasure \\
1 & krh (كره) & 41 & 14 & 27 & Medinan-tilted (\emph{p} = 0.005);
twofold \emph{karāha}/\emph{ikrāh} \\
2 & ḍyq (ضيق) & 13 & 9 & 4 & Inner constriction (\emph{ḍīq al-ṣadr};
\emph{ḍāqa dharʿan}) \\
2 & ḥzn (حزن) & 42 & 26 & 16 & Most frequent Stage-2 lexeme; valence
transversal (see §4.3.2) \\
2 & ʾsf (أسف) & 5 & 5 & 0 & Exclusively Meccan \\
3 & nqm (نقم) & 17 & 12 & 5 & Meccan-tilted; vengeful disapproval;
cf.~\emph{al-muntaqim} \\
3 & sxṭ (سخط) & 4 & 0 & 4 & Exclusively Medinan (\emph{p} = 0.005); see
baseline-sensitivity analysis in §4.5.2 \\
3 & mqt (مقت) & 6 & 4 & 2 & Most intense moral aversion (Zajjāj) \\
4 & ġḍb (غضب) & 24 & 13 & 11 & Central anger lexeme; predicated of God,
prophets, believers \\
4 & ḥrd (حرد) & 1 & 1 & 0 & Qurʾānic hapax: Q. 68:25 \\
5 & ġyẓ (غيظ) & 11 & 3 & 8 & Compressed anger; \emph{tamayyuz}
manifestation (Q. 67:8) \\
6 & bġy (بغي) & 96 & 55 & 41 & Bridge node; bidirectional causation
(§4.6 caveats) \\
6 & ṭġy (طغي) & 39 & 29 & 10 & Pharaonic; structural arrogance \\
6 & ʿtw (عتو) & 10 & 9 & 1 & Settled defiance; properly
\emph{istikbār}-attached \\
\textbf{Σ} & --- & \textbf{312} & \textbf{183} & \textbf{129} & \\
\end{longtable}

Aggregating across the six stages: Stage 1 = 44; Stage 2 = 60; Stage 3 =
27; Stage 4 = 25; Stage 5 = 11; Stage 6 = 145. The asymptotic
chi-squared test against the null of uniform six-stage distribution
rejects (χ²(5) = 227.15, \emph{p} \textless{} 0.001), with effect sizes
Cramér's \emph{V} = 0.38 and Cohen's \emph{w} = 0.85 (large effect by
Cohen's conventions). However, an asymptotic χ² on a six-cell
distribution with one cell at \emph{n} = 11 (Stage 5) and another at
\emph{n} = 145 (Stage 6) is not above suspicion: the chi-squared
approximation requires expected counts ≥ 5 in every cell (which is met)
but is \emph{known} to be conservative under marginal imbalance (Cochran
1954). We therefore complement the asymptotic test with a
\textbf{marginal-preserving permutation null}: the per-root attestation
counts are held fixed and the fourteen stage labels are randomly
permuted across roots (10,000 permutations, fixed RNG seed 20260509),
yielding the empirical \emph{p}-value as the fraction of permuted χ²
values that meet or exceed the observed 227.15. The result is \emph{p} =
0.24 --- i.e.~under random stage relabelling of the fourteen roots, a χ²
statistic this extreme is not unusual, because much of the omnibus's
mass derives from the fact that \emph{baghy} alone contributes 96
attestations (versus a uniform expectation of \textasciitilde52 per
stage) and would inflate any single-stage cell into which it fell. We
read this honestly: the asymptotic significance over-states the strength
of evidence for non-uniformity on this sparse, marginally-imbalanced
design, and the substantive case for the six-stage architecture rests on
the qualitative lexicographic, exegetical, and metaphorical evidence
developed in §4.1 and §4.3--§4.8 rather than on the omnibus
\emph{p}-value alone. We retain the asymptotic χ² as a descriptive
statistic with appropriate effect-size reporting, but treat it as
\emph{consistent with} --- not as \emph{positive evidence for} --- H1. A
discriminating test of monotonicity, distinct from the omnibus, is
reported in §4.2.1 below. A complementary chi-squared test against the
null of equal split between the phenomenological core (Stages 1--5 = 167
attestations) and the behavioural-outcomes pole (Stage 6 = 145
attestations) \emph{fails to reject} (χ²(1) = 1.55, Cohen's \emph{w} =
0.07 --- trivial effect). Because failure to reject a null is
\emph{absence of evidence of a difference} and not evidence of equality,
we report this result not as positive corroboration of H2 but as a
finding \emph{consistent with} the reading developed in §4.8.5 and §5.4:
the Stage-6 lexemes (\emph{baghy}, \emph{ṭughyān}, \emph{ʿutuww}) are
not exclusively anger-derived but stand in a \emph{bidirectional} causal
relation to the anger spectrum---anger may produce these behaviours, but
the behaviours can equally produce anger (in their victims, the
prophets, and the divine response). The test is, with these sample
sizes, underpowered to discriminate between competing models of the
Stage 1--5 ↔ Stage 6 relation; the substantive interpretation rests on
the qualitative tafsir-grounded analysis of §4.8.5, not on the
statistical near-balance.

\paragraph{4.2.1. Monotonic-trend tests for the six-stage
ordering}\label{monotonic-trend-tests-for-the-six-stage-ordering}

Because the χ²(5) test confirms non-uniformity but says nothing about
\emph{ordering}, we report three complementary monotonic-trend tests of
stage rank against attestation frequency. (i) \textbf{Spearman
rank-correlation} between stage index (1--6) and per-stage attestation
count is ρ = --0.09 with two-sided asymptotic \emph{p} = 0.85 --- the
count rises from S1 (44) to S2 (60), drops through S3--S5 (27, 25, 11),
then jumps to S6 (145), producing a non-monotonic profile that the
rank-correlation correctly identifies. (ii) \textbf{Mann-Kendall}
(small-sample-robust) yields \emph{S} = --3, \emph{z} = --0.38, \emph{p}
= 0.71. (iii) \textbf{Cochran-Armitage-style} OLS slope of count on
stage rank yields slope = 10.17, \emph{z} = 0.85, \emph{p} = 0.40. All
three tests fail to reject the no-trend null. We read this as the
expected and substantively interpretable result: the action-intensity
continuum is an ordering of \emph{semantic intensity}, not of
\emph{attestation frequency}; the Qurʾān's pedagogical emphasis at the
lower-intensity (S1--2) and behavioural-outcome (S6) poles is
theoretically motivated rather than indexical of intensity. The
non-monotonic frequency profile is precisely what one would expect if
the corpus weights pre-anger and post-anger lexemes pedagogically,
\emph{because} of their hortatory function in the moral architecture of
the text. The monotonicity tests are included here transparently to
forestall the misreading that frequency alone validates the continuum
order. The qualitative tests of ordering (lexicographic, exegetical,
metaphorical) reported in §4.1 and §4.3--§4.8 are the load-bearing
evidence for the continuum's structure.

\begin{figure}
\centering
\includegraphics[width=0.95\linewidth,height=\textheight,keepaspectratio]{fig2_frequency_by_stage.pdf}
\caption{Per-root and per-stage frequency of the 312 core spectrum
attestations. \emph{Left panel}: linear-scale frequency for each of the
fourteen roots; bar colour encodes the phenomenological stage. Stage 6
(\emph{bġy} 96 vs.~\emph{ʿtw} 10) shows a strong internal asymmetry that
mirrors the lexical division of labour between distributed-aggression
(\emph{baghy}) and obstinate-defiance (\emph{ʿutuww}) outlined in §4.6.
\emph{Right panel}: stage totals (S1=44, S2=60, S3=27, S4=25, S5=11,
S6=145). The six-stage distribution rejects uniformity at χ²(5) =
227.15, \emph{p} \textless{} 0.001. The phenomenological core (S1--5 =
167) and the behavioural-outcomes pole (S6 = 145) are \emph{not}
significantly unequal in raw count (χ²(1) = 1.55, fails to reject), a
result we read as supporting the bidirectional-causation analysis of
\emph{baghy}/\emph{ṭughyān}/\emph{ʿutuww} developed in §4.6 and §5.4 ---
the Qurʾān is not pedagogically over-loading anger-outcomes; rather,
Stage-6 lexemes are themselves not exclusively anger-derived.}
\end{figure}

\begin{figure}
\centering
\includegraphics[width=0.95\linewidth,height=\textheight,keepaspectratio]{fig3_meccan_medinan.pdf}
\caption{Stacked Meccan/Medinan distribution for each of the fourteen
core roots. The numerals above each bar give the raw Meccan/Medinan
split. Two distributional extremes warrant emphasis. \emph{ʾasaf} (Stage
2) is attested \textbf{exclusively} in Meccan suras (5/0): the lexeme of
``regret-grief'' emerges in the discursive context of prophetic
disappointment with rejection (Q. 18:6, Q. 43:55, Q. 7:150).
\emph{Sakhaṭ} (Stage 3) is attested \textbf{exclusively} in Medinan
suras (0/4; binomial \emph{p} = 0.005 against the corpus-derived 0.74
baseline of Meccan ayat): the lexeme of ``comprehensive divine
displeasure'' emerges only after the Hijra, in confrontation with the
\emph{munāfiqūn} (Q. 47:28, Q. 5:80, Q. 9:58, Q. 3:162). Additional
patterns: \emph{naqm} (Stage 3) is Meccan-tilted (12/5); \emph{maqt}
(Stage 3) is balanced (4/2); \emph{ḥard} is a Qurʾānic hapax at Q. 68:25
(Meccan). The visualisation thus encodes a non-trivial finding: not
every root in the spectrum behaves as a context-neutral lexical resource
--- several of them are \emph{discursive-contextually specialised}, a
pattern theorised in §5.5 as a general feature of Qurʾānic vocabulary.}
\end{figure}

\begin{figure}
\centering
\includegraphics[width=0.95\linewidth,height=\textheight,keepaspectratio]{fig5_morphology_stack.pdf}
\caption{Morphological profile of the fourteen core roots, broken out by
part of speech. The left-to-right gradient shows a \emph{typological
shift} across the continuum that supplements the frequency analysis.
Stages 1--2 are verbally dominant --- \emph{ḥzn} 88\% verb forms (37 v /
5 n), \emph{ḍyq} 62\% --- consistent with the ontology of pre-anger
inner states as \emph{transient, situationally-induced} phenomena. Stage
4 nominalises: \emph{ġḍb} shifts the centre of gravity to the noun (16 n
/ 6 v), as the lexeme acquires a stable, attributable referent (``the
wrath of God''). Stage 6 is the most nominal, with participial forms
emerging (\emph{ṭāgūt}, \emph{bāghin}); the behavioural-outcomes
vocabulary entrenches as a stable agent-property rather than a transient
process. This morphological gradient --- verbal Stages 1--2 → nominal
Stage 6 --- is independent corpus-internal evidence for the
action-intensity continuum: as intensity ascends, the linguistic
representation reifies from \emph{event} to \emph{attribute} to
\emph{identity}.}
\end{figure}

\subsubsection{4.3. Stage 1 --- Pre-anger
displeasure}\label{stage-1-pre-anger-displeasure}

The lower bound of the continuum is occupied by two roots---\emph{ʾ-f-f}
and \emph{k-r-h}---that mark the \emph{pre-anger} register: vocal
irritation and inner aversion respectively.

\textbf{4.3.0. \emph{Uff} (3 attestations; all Meccan).} \emph{Uff}
(gloss: \emph{taḍajjur}, ``irritation'') is the lowest possible verbal
expression of displeasure in the Qurʾānic register. Its three
attestations occupy paradigmatic moral and rhetorical positions:

\begin{itemize}
\tightlist
\item
  Q. 17:23 prohibits saying \emph{uff} to one's parents (\emph{fa-lā
  taqul lahumā uffin}) --- the \emph{lowest} possible verbal sign of
  displeasure is morally weighted at the highest possible relational
  stake.
\item
  Q. 21:67 places \emph{uff} in Abraham's mouth as he confronts his
  idol-worshipping people (\emph{uffin lakum wa-li-mā taʿbudūna min dūni
  Allāh}) --- a prophetic remonstrance.
\item
  Q. 46:17 ascribes \emph{uff} to a rebellious child reproaching his
  believing parents --- the inverse of the Q. 17:23 case.
\end{itemize}

As the lowest tier of the spectrum, \emph{uff} establishes a
foundational point of the Qurʾānic discourse: even the slightest verbal
sign of displeasure is morally weighted. The root therefore anchors the
lower boundary of the anger spectrum at a register that is below
``anger'' proper but already morally accountable.

\textbf{4.3.1. \emph{Karh} (41 attestations; 14 Meccan / 27 Medinan;
binomial \emph{p} = 0.005 in the Medinan direction).} \emph{Karh}
exhibits a \emph{twofold structure} that must be analytically
distinguished:

\begin{enumerate}
\def\labelenumi{(\alph{enumi})}
\item
  \textbf{Inner \emph{karāha} (dislike).} As an affective state,
  \emph{karāha} is the agent's pre-anger evaluative dislike of an
  object, person, or course of action. Q. 2:216 \emph{ʿasā an takrahū
  shayʾan wa-huwa khayrun lakum} (``perhaps you dislike a thing while it
  is good for you'') frames inner dislike as a \emph{reformable}
  affective state --- its truth-value about the object can be revised by
  knowledge or by divine corrective. This is properly the Stage-1
  \emph{anger-spectrum} sense.
\item
  \textbf{Structural \emph{ikrāh} (coercion).} \emph{Ikrāh} is
  \emph{not} an emotion but a coercion mechanism: the structural
  extension of one party's distaste into compulsion of another. Q. 2:256
  \emph{lā ikrāha fī al-dīn} (``there is no coercion in religion'') sets
  the ethical limit on this structural extension; Q. 16:106 (\emph{illā
  man ukriha}, ``except one who is forced'') provides the legal
  exemption for testimony given under coercion; Q. 24:33 (\emph{lā
  tukrihū fatayātikum ʿalā al-bighāʾ}) prohibits the coercion of
  slave-women into prostitution. We include \emph{karh} in the spectrum
  on the strength of its \emph{karāha} sense; the \emph{ikrāh} sense
  lies outside the affective field but is a structural cognate worth
  noting because it shares the same root and instantiates the
  \emph{behavioural extension} of distaste --- itself a Stage-6
  phenomenon ahead of its time.
\end{enumerate}

\subsubsection{4.4. Stage 2 --- Inner pressure and
contraction}\label{stage-2-inner-pressure-and-contraction}

Stage 2 brings together three roots---\emph{ḍ-y-q}, \emph{ḥ-z-n},
\emph{ʾ-s-f}---with 60 combined attestations. Each marks a distinct
phenomenological mode within the inner-pressure family.

\textbf{4.4.1. \emph{Ḍīq} (13; 9M/4Md).} Ibn Fāris glosses \emph{ḍyq} as
``constriction, narrowness''; al-Rāghib in opposition to \emph{saʿah}
(expansion). The root appears predominantly in the construction
\emph{ḍīq al-ṣadr} (constriction of the chest), localising the
experience to the embodied site. Q. 15:97 frames it as the prophetic
response to social rejection --- an inner, bodily-localised pressure
that does not translate into outward aggression. Al-Ṭabāṭabāʾī
(\emph{al-Mīzān} 12:195) observes the verse expresses divine
\emph{knowledge} of prophetic \emph{ḍīq} rather than prohibition: a
phenomenologically natural, theologically licit state, not a moral
failing. Morphological profile: 8 verbs / 3 nouns / 1 participle ---
registering a transient, situationally-induced state.

\textbf{4.4.1.1. \emph{Ḍāqa dharʿan} --- inhibited aggression at Lot, Q.
11:77/80.} The compound idiom \emph{ḍāqa bihim dharʿan} (``his arm-span
tightened by them'') in the Lot narrative figures \emph{post-anger
blocked aggression} --- anger at action-threshold but structurally
inhibited. The four-step trajectory (trigger \emph{sīʾa bihim} →
constriction \emph{ḍyq+dharʿ} → aspirational voicing \emph{law anna lī
bikum quwwatan} → transfer of agency to the divine through the angels'
arrival) shows that \emph{ḍyq} sits at \emph{both} poles of the
spectrum: pre-anger contraction \emph{and} upper-pressure threshold of
inhibition. The Qurʾānic narrative theologises inhibited aggression:
when the wronged agent lacks power, divine agency completes the anger
trajectory.

\textbf{4.4.2. \emph{Ḥuzn} (42; 26M/16Md).} The most frequent and most
verbal Stage-2 lexeme (37v / 5n = 88\% verbal). Ibn Manẓūr preserves the
etymologically primary ``rugged, uneven ground'' (\emph{al-arḍ
al-ġalīẓah}); the metaphorical extension figures grief as inner
roughness, anticipating Kövecses's (2000) EMOTION IS A LANDSCAPE schema.
Q. 2:38 places \emph{ḥuzn} in syntagmatic pair with \emph{khawf} and
opposes both to divinely-given guidance --- \emph{ḥuzn} as the affective
signature of \emph{separation from guidance}; al-Ṭabāṭabāʾī
(\emph{al-Mīzān} 1:146) calls this foundational for \emph{ʿilm al-nafs
al-Qurʾānī}. \textbf{Caveat (transversality)}: \emph{Ḥuzn} is \emph{not}
strictly anger-derived. It has four valences --- standalone grief over
loss (Q. 12:84, Jacob over Joseph), externally-induced sorrow from being
wronged (prophetic grief at rejection), positive moral compunction (the
affective register of \emph{tawba}), and anger turned inward when
externally inhibited (closest to the Stage-2 placement here). Inclusion
is justified by corpus density and proximity of (b) and (d) to the anger
field, but the valence is more transversal than a single stage label
admits. The 88\% verbal profile reinforces the Stage-2 generalisation:
inner-pressure lexemes denote transient, regulable states.

\textbf{4.4.3. \emph{Asaf} (5; 5M/0Md).} The least frequent Stage-2
lexeme. Ibn Fāris and al-Rāghib gloss as the conjunction of grief and
anger (\emph{al-jamʿ bayna al-ḥuzn wa-l-ġaḍab}) --- at the threshold
between inner pressure and active anger. Q. 43:55 (\emph{fa-lammā
āsafūnā intaqamnā minhum}) frames \emph{asaf} as the immediate causal
precursor of divine retribution; al-Ṭabarsī (\emph{Majmaʿ} 9:82) glosses
\emph{al-asaf} here as ``the intense anger arising from prior grief.''
The exclusive-Meccan distribution (5/0; raw \emph{p} = 0.34, adj.
\emph{p} = 1.00) is suggestive but evidentially weak --- at \emph{n} = 5
the test has very low power. Across all five attestations, \emph{asaf}
occurs in narrative passages depicting prior prophets' encounter with
rejection and in the \emph{ghaḍbān asifan} compound applied to Moses
(§4.4.3.1). We register the 5/0 distribution as a \emph{contextual
observation} warranting follow-up at higher \emph{n} (e.g.~against the
\emph{ḥadīth} corpus), not as a confirmed empirical finding.

\textbf{4.4.3.1. \emph{Ghaḍbān asifan} --- sacred anger at Moses, Q.
7:150 / 20:86.} The compound \emph{ghaḍbān asifan}, applied to Moses on
his return to find the calf-worshippers, generates a three-layer
dialectic: \emph{ghaḍab} is outward/kinetic (throwing the tablets,
seizing Aaron); \emph{asaf} is inward/potential (compassionate grief
over the relapse). \emph{Asaf} grounds \emph{ghaḍab}: absent its
compassionate ground, Moses' anger would be violence; present, it is
anger in service of guidance. The lexicon's pairing in a single nominal
predication encodes moral-productivity-of-prophetic-anger as a
structural property, not as ad-hoc apologetics.

\textbf{4.4.4. Network evidence for Stage-2 coherence.} The aya-level
co-occurrence network (Figure 5) reveals that the strongest single edge
connects \emph{ḥzn} and \emph{ḍyq} (weight 3); these two lexemes share
three verses in the Qurʾān, more than any other pair within the
spectrum. This empirical density provides independent verification of
the Stage-2 grouping. Centrality analysis (Figure 6) places \emph{ḥzn}
among the highest-degree nodes (weighted degree 5, tied with
\emph{ġḍb}).

\subsubsection{4.5. Stage 3 --- Evaluative
aversion}\label{stage-3-evaluative-aversion}

Stage 3 brings together three roots---\emph{n-q-m}, \emph{s-kh-ṭ},
\emph{m-q-t}---with 27 combined attestations. Each marks a distinct mode
of \emph{evaluative-aversive} response to a moral object.

\textbf{4.5.1. \emph{Naqm} (17 attestations; 12 Meccan / 5 Medinan).}
\emph{Naqm} is the lexeme of \emph{vengeful disapproval}: anger that is
grounded in moral evaluation and oriented toward retribution. Q. 7:126
(\emph{wa-mā naqamū minnā illā an āmannā bi-āyāti rabbinā}) is the
linguistic exemplar: anger directed \emph{at others' belief} ---
Pharaoh's sorcerers, on accepting Moses' message, are described as
suffering Pharaonic \emph{naqm} precisely \emph{because} they believed.
The grammatical construction \emph{naqamū min} + (cause) is a canonical
Qurʾānic formula for evaluatively-grounded anger.

The lexeme also generates \emph{al-muntaqim} (``the Avenger'') as a
divine attribute (Q. 32:22 and parallels), establishing a direct
theological-lexical bridge from the human-affective register to the
divine retributive register. Position on the spectrum: \emph{naqm} sits
between \emph{sakhaṭ} and \emph{ghaḍab} on the action axis. It is more
action-oriented than \emph{sakhaṭ} (which is a stable moral
disapprobation) but less kinetic than \emph{ghaḍab} (which mobilises the
motor agent). The Meccan tilt (12/5) is consistent with the lexeme's
frequent appearance in narrative-prophetic passages (the Pharaonic
sorcerers, the people of Lot, the people of Madyan).

\textbf{4.5.2. \emph{Sakhaṭ} (4 attestations; 0 Meccan / 4 Medinan).}
Despite its low frequency, \emph{sakhaṭ} furnishes one of the paper's
principal distributional findings. The exclusive-Medinan
distribution---zero Meccan attestations among four total---has a
two-sided exact binomial p-value of \textbf{0.005} against the
corpus-derived baseline of 0.74 Meccan probability, statistically
significant at α = 0.05.

\textbf{\emph{Baseline derivation, made explicit.}} The 0.74 Meccan-rate
baseline is the proportion of \emph{ayat} (verses, not suras) classified
Meccan in the QAC sura metadata, computed as
\texttt{meccan\_ayat\ ÷\ total\_ayat} weighted by \texttt{total\_verses}
per sura. We use the ayat-weighted rate rather than the sura-count rate
(\textasciitilde0.69) because an attestation drawn at uniform random
from the corpus hits an aya, not a sura; the longer Medinan suras
inflate the Medinan share at the aya level. The implementation is in
\texttt{analysis/advanced\_metrics.py}
(\texttt{\_compute\_meccan\_baseline}).

\textbf{\emph{Granular baseline sensitivity.}} With n = 4, the inference
is mildly sensitive to the baseline. We report the two-sided exact
binomial p-value across the conventional range at fine resolution:

\begin{longtable}[]{@{}
  >{\centering\arraybackslash}p{(\linewidth - 4\tabcolsep) * \real{0.3333}}
  >{\centering\arraybackslash}p{(\linewidth - 4\tabcolsep) * \real{0.3333}}
  >{\raggedright\arraybackslash}p{(\linewidth - 4\tabcolsep) * \real{0.3333}}@{}}
\toprule\noalign{}
\begin{minipage}[b]{\linewidth}\centering
H₀ Meccan rate
\end{minipage} & \begin{minipage}[b]{\linewidth}\centering
Two-sided p (k=0, n=4)
\end{minipage} & \begin{minipage}[b]{\linewidth}\raggedright
Interpretation
\end{minipage} \\
\midrule\noalign{}
\endhead
\bottomrule\noalign{}
\endlastfoot
0.70 (lower bound; sura-count proxy) & 0.0081 & significant at α =
0.05 \\
0.72 & 0.0061 & significant at α = 0.05 \\
\textbf{0.74 (corpus ayat-weighted; default)} & \textbf{0.0046} &
\textbf{significant at α = 0.05} \\
0.76 & 0.0033 & significant at α = 0.01 \\
0.78 (upper bound) & 0.0023 & significant at α = 0.01 \\
\end{longtable}

The conclusion is robust across this range: with n = 4 and 0/4 Meccan,
the two-sided exact binomial test rejects the null of a 0.70--0.78
Meccan rate at α = 0.05.

\textbf{\emph{Multiple-comparison adjustment.}} Step-down
Holm-Bonferroni over the family of fourteen per-root binomials yields
\textbf{adjusted \emph{p} = 0.05 (raw 0.005)} for \emph{sakhaṭ} ---
borderline rather than emphatic significance. \emph{Asaf} (5/0; raw
\emph{p} = 0.34, adj. 1.00), \emph{ʿutuww} (9/1; raw 0.47, adj. 1.00),
and \emph{karh} (raw 0.005, adj. 0.05 --- tied with \emph{sakhaṭ}) tilts
do not survive correction. With \emph{n} = 4 the test is underpowered
against subtler alternatives; the substantive claim of discursive
specialisation rests on the convergence of distributional evidence with
the contextual-tafsir reading, not on the binomial alone.

Contextual analysis of all four occurrences reveals the same operative
pattern: \emph{sakhaṭ} attests in connection with the \emph{Munāfiqūn}
(hypocrites) or with quasi-believers whose outward profession of faith
conflicts with their inward dispositions. Q. 47:28 (\emph{ittabaʿū mā
askhaṭa Allāh wa-karihū riḍwānahu}) is the canonical exemplar. Following
Narimani et al.~(2021), we read \emph{sakhaṭ} not as an instantaneous
emotional surge but as a sustained \emph{ethical disapprobation}; the
present paper extends their lexical-tafsīr argument with corpus evidence
and proposes:

\begin{quote}
\textbf{Hypothesis (Discursive Specialisation of \emph{sakhaṭ}):} In the
Qurʾānic lexicon, \emph{sakhaṭ} operates as a discursively specialised
lexeme, reserved for divine displeasure with the
\emph{Munāfiqūn}---agents whose \emph{zāhir} (outward profession)
departs from their \emph{bāṭin} (inward orientation). Because the
\emph{Munāfiqūn} category is constitutively a Medinan-period phenomenon,
\emph{sakhaṭ} attests exclusively in Medinan suras. This pattern is an
instance of a more general phenomenon of \emph{discursive-contextual
lexical specialisation} in the Qurʾānic register.
\end{quote}

\textbf{4.5.3. \emph{Maqt} (6 attestations; 4 Meccan / 2 Medinan).}
\emph{Maqt} is \emph{intense moral aversion}: the four lexica converge
on \emph{ashaddu al-ibghāḍ} (Zajjāj/al-Layth in \emph{Tahdhīb
al-Lugha}); al-Rāghib grounds it specifically in observation of a
reprehensible act (\emph{ʿan amrin qabīḥin rakibahu}), making
\emph{maqt} purely evaluative. Two contrasts with \emph{ghaḍab}:
causation (purely evaluative vs.~injury/frustration-triggered); action
profile (rejection-distance vs.~approach-aggression). Canonical
exemplars --- Q. 40:10 (\emph{maqt al-nafs} doubled), Q. 4:22
(institutional moral aversion in forbidden marriages), Q. 35:39 (rising
divine \emph{maqt} against disbelievers), Q. 61:3 (saying-without-doing)
--- stabilise the lexeme as the Qurʾānic register's term for moral
aversion as evaluative response, distinct from the kinetic-motor
\emph{ghaḍab} of Stage 4.

\subsubsection{4.6. Stage 4 --- Active
anger}\label{stage-4-active-anger}

Stage 4 (25 attestations) covers active, motor-mobilising anger.

\textbf{4.6.1. \emph{Ghaḍab} (24; 13M/11M).} The central anger lexeme of
the Qurʾān, with near-balanced Meccan/Medinan distribution marking it as
the unmarked default. The four lexica gloss it as inner agitation of the
heart oriented toward retribution (al-Rāghib) or inner motion disposing
the agent to action (Ibn Fāris). Predicated balance-fully of God (Q.
48:6), prophets, and believers. Morphologically nominal (16 n / 6 v / 2
ptcp = 67\% nominal vs \emph{ḥuzn}'s 88\% verbal): \emph{ghaḍab} is a
\emph{structural concept-state} with organisational ontological standing
rather than a transient affect. Al-Fakhr al-Rāzī's \emph{Mafātīḥ
al-Ghayb} glosses divine \emph{ghaḍab} as \emph{irādat inzāl al-ʿiqāb}
--- a normative-juridical orientation, not an emotional response.
Network analysis places \emph{ghḍb} among the top three for betweenness
centrality (point estimate 0.17), with notable edges to \emph{ġyẓ}
(upward transition to S5), \emph{ḥzn} (downward S2 link), and \emph{bġy}
(S6 link).

\textbf{4.6.3. \emph{Ḥard} (Q. 68:25, hapax).} A Qurʾānic hapax in the
Garden-Owners parable (\emph{wa-ghadaw ʿalā ḥardin qādirīn}):
garden-owners who vow in the night to harvest before dawn so as to deny
the poor their share. Classical exegesis converges on a dual gloss:
al-Ṭabarsī's \emph{Majmaʿ al-Bayān} reads \emph{ḥard} as both
\emph{ghaḍab} (anger) and \emph{qaṣd} (purposive intent);
al-Ṭabāṭabāʾī's \emph{al-Mīzān} refines as anger-plus-deliberate-denial
--- purposive anger weaponised through resource-denial. \emph{Ghaḍab}
mobilises the motor agent outward; \emph{ḥard} mobilises toward
deliberate, miserly withholding --- anger that strips the world of what
it is owed. The narrative consequence is the destruction of the garden
by night-storm: divine retribution against the \emph{ḥard}-driven plan
to deny. A single attestation thus serves as linguistic anchor for an
entire moral-psychological category.

\subsubsection{4.7. Stage 5 --- Compressed/explosive
rage}\label{stage-5-compressedexplosive-rage}

\textbf{4.7.1. \emph{Ghayẓ} and the \emph{kaẓm} metaphor (11
attestations; 3M/8Md).} The four lexica converge: \emph{ghayẓ} is
``compressed contained anger'' (al-Rāghib), ``the highest concentration
of restrained anger'' (Mostafavi), ``filling of the self with restrained
anger'' (Ibn Manẓūr) --- pressurised contents of a sealed vessel. Q.
3:134 commends \emph{al-kāẓimīn al-ghayẓ}; the verb \emph{kaẓm} derives
from the act of tying off a waterskin, the agent-internal regulation
that sustains containment. Read through Lakoff \& Kövecses (1987), the
\emph{ghayẓ--kaẓm} pair instantiates ANGER IS A PRESSURIZED CONTAINER
with precision: \emph{ghayẓ} the contained fluid, \emph{kaẓm} the seal
(Figure 7). The Medinan tilt (8/3) tracks the discursive context: inner
anger management becomes a salient pedagogical concern in the formative
Medinan community facing internal \emph{Munāfiq} and external Quraysh
hostility.

\textbf{4.7.2. \emph{Tamayyuz}: container collapse in Q. 67:8
(manifestation, not separate root).} Q. 67:8 --- \emph{takādu tamayyazu
min al-ġayẓ} (``{[}the fire{]} is on the verge of bursting apart from
rage'') --- is the single most articulated image of rage in the Qurʾān.
\emph{Tamayyaza} (root \emph{m-y-z}: to part, sever) figures not the
\emph{exit} of pressurised content (as in English \emph{she blew her
top}, \emph{he exploded}, \emph{I lost it}) but the \emph{structural
integrity of the container itself failing}. We treat \emph{tamayyuz} as
a manifestation of Stage 5 rather than a distinct core root, retaining
the fourteen-root inventory. H5 is corroborated in the form: Q. 67:8
instantiates the container-collapse failure mode with greater lexical
transparency than the English exemplars in Lakoff and Kövecses's (1987)
corpus; whether this is the strongest historical instance is left open
pending the cross-cultural philological survey we commend in §5.6.

\subsubsection{4.8. Stage 6 --- Behavioural outcomes (with
caveats)}\label{stage-6-behavioural-outcomes-with-caveats}

\textbf{4.8.1. \emph{Baghy} (96; 55M/41Md).} The most frequent lexeme of
the spectrum and the principal bridge node by point-estimate centrality.
Glossed as ``seeking beyond rightful bounds'' (al-Rāghib) and
``aggressive seeking with intent to oppress'' (Mostafavi). Q. 42:42
collocates \emph{baghy} with \emph{ẓulm}, embedding it in the
moral-evaluation register. Ibn ʿĀshūr clarifies (\emph{al-Taḥrīr}
25:126): \emph{baghy} and \emph{ẓulm} are near-synonymous, except
\emph{baghy} is specific to non-rightful aggression against another
while \emph{ẓulm} covers all injustice.

\textbf{4.8.2. \emph{Baghy} as bridge node (with bootstrap
uncertainty).} \emph{Baghy} attains the highest point-estimate closeness
centrality (0.55) and is tied with \emph{ghaḍab} (0.53) and \emph{asaf}
(0.49) for the highest betweenness centrality (0.15). Bootstrap 95\%
percentile CIs (1,000 ayat-resamples) are wide --- for \emph{baghy}:
betweenness ∈ {[}0.00, 0.31{]}; closeness ∈ {[}0.23, 0.84{]};
eigenvector ∈ {[}0.00, 0.68{]} --- reflecting the small graph (14 nodes)
and modest edge density. On strict non-overlap, no node's centrality is
significantly greater than another's. The bridging-role claim must
therefore be qualified: \emph{baghy} leads on three of four
point-estimate measures, but the bootstrap does not rule out
alternatives. What is \emph{robust} at the aggregate level is the
cross-field connectivity to the umbrella moral terms (Table 4):
\emph{baghy} shows 16 aya-level co-occurrences with
\emph{ẓlm/jrm/fsq/ʿdw/šnʾ} combined (8 with \emph{ʿdw}, 4 with
\emph{ẓlm}), an integer-count pattern independent of centrality ranking.
H4 is \emph{consistent with} point estimates, qualified by bootstrap.

\textbf{Table 4. Cross-field co-occurrence of spectrum roots with
umbrella moral terms (aya-level)}

Cross-field ties for the original ten-root analysis (preserved here from
the prior pipeline run; ties for the four newly-incorporated roots ---
\emph{ʾff}, \emph{naqm}, \emph{maqt}, \emph{ḥard} --- are sparse and
dominated by zeros, given their low corpus frequency, and are omitted
from this summary table without affecting the overall pattern):

\begin{longtable}[]{@{}
  >{\centering\arraybackslash}p{(\linewidth - 12\tabcolsep) * \real{0.1429}}
  >{\centering\arraybackslash}p{(\linewidth - 12\tabcolsep) * \real{0.1429}}
  >{\centering\arraybackslash}p{(\linewidth - 12\tabcolsep) * \real{0.1429}}
  >{\centering\arraybackslash}p{(\linewidth - 12\tabcolsep) * \real{0.1429}}
  >{\centering\arraybackslash}p{(\linewidth - 12\tabcolsep) * \real{0.1429}}
  >{\centering\arraybackslash}p{(\linewidth - 12\tabcolsep) * \real{0.1429}}
  >{\centering\arraybackslash}p{(\linewidth - 12\tabcolsep) * \real{0.1429}}@{}}
\toprule\noalign{}
\begin{minipage}[b]{\linewidth}\centering
Root
\end{minipage} & \begin{minipage}[b]{\linewidth}\centering
šnʾ
\end{minipage} & \begin{minipage}[b]{\linewidth}\centering
jrm
\end{minipage} & \begin{minipage}[b]{\linewidth}\centering
fsq
\end{minipage} & \begin{minipage}[b]{\linewidth}\centering
ẓlm
\end{minipage} & \begin{minipage}[b]{\linewidth}\centering
ʿdw
\end{minipage} & \begin{minipage}[b]{\linewidth}\centering
Total
\end{minipage} \\
\midrule\noalign{}
\endhead
\bottomrule\noalign{}
\endlastfoot
ḍyq & 0 & 0 & 0 & 0 & 0 & 0 \\
ḥzn & 0 & 0 & 0 & 1 & 1 & 2 \\
ʾsf & 0 & 0 & 0 & 1 & 1 & 2 \\
sxṭ & 0 & 0 & 0 & 0 & 0 & 0 (krh-tie excluded since \emph{krh} is now in
the spectrum) \\
ġḍb & 0 & 0 & 0 & 2 & 3 & 5 \\
ġyẓ & 0 & 0 & 0 & 0 & 1 & 1 \\
\textbf{bġy} & \textbf{1} & \textbf{1} & \textbf{2} & \textbf{4} &
\textbf{8} & \textbf{16} \\
ṭġy & 0 & 0 & 0 & 2 & 1 & 3 \\
ʿtw & 0 & 0 & 0 & 0 & 0 & 0 \\
\end{longtable}

The radical concentration of cross-field ties at \emph{baghy} is
striking. Of the total cross-field ties across the fourteen spectrum
roots, 19 (\textasciitilde53\%) involve \emph{baghy}. \emph{Baghy} thus
serves as the principal lexical bridge between the emotional and the
broader moral-evaluation field of the Qurʾān.

\textbf{4.8.3. \emph{Ṭughyān} (39 attestations; 29 Meccan / 10
Medinan).} Root \emph{ṭ-gh-y} derives from water flooding past its
banks. Q. 96:6--7 (\emph{kallā inna l-insāna la-yaṭghā an raʾāhu
staġnā}) compresses a psychology of rebellion into a conditional:
\emph{ṭughyān} arises from the appraisal of \emph{istighnāʾ}
(self-sufficiency) --- structurally an appraisal-theoretic claim,
congruent with Lazarus (1991) and Scherer (2001, 2005), and aligning
with Roseman's ``agency-self / motive-inconsistent / high-control''
pattern. Qarizadeh et al.'s (2024) polysemy analysis notes the cognitive
structure; we identify \emph{ṭughyān} as the cognitive-appraisal node of
Stage 6.

\textbf{4.8.4. \emph{ʿUtuww} (10 attestations; 9 Meccan / 1 Medinan).}
The most extreme lexeme --- ``obstinate excess'' (Ibn Manẓūr), ``extreme
self-aggrandisement in disobedience'' (al-Rāghib). Q. 25:21 collocates
\emph{ʿatū ʿutuwwan kabīrā} with \emph{istakbarū fī anfusihim}:
rebellion crystallised into existential posture. Strong Meccan
distribution (9/1; raw \emph{p} = 0.06; Holm-adjusted \emph{p} = 1.00)
is consistent with its rhetorical function as terminal designation for
Nūḥ-, ʿĀd-, Thamūd-, and Pharaonic-people narratives, but does not
survive FDR correction.

\textbf{4.8.5. Bidirectional causation and a tafsir-grounded audit.}
\emph{Baghy} arises in Qurʾānic usage frequently from greed and
\emph{ʿulw fī al-arḍ} (cf.~Q. 6:21) rather than from acute anger;
\emph{ṭughyān} --- paradigmatically Pharaonic (Q. 79:17) --- is rooted
in structural arrogance, with Q. 26:55 (\emph{innahum lanā la-ghāʾiẓūn})
reversing the polarity to show oppressor-anger arising \emph{from} the
rebellion of the oppressed; \emph{ʿutuww} attaches to \emph{istikbār} as
settled character. We retain Stage 6 because the corpus locates these at
the behavioural pole, but the causation arrow is \emph{bidirectional}.

To operationalise this caveat, each of the 145 Stage-6 attestations was
coded on a tafsir-grounded ternary against the joint reading of
\emph{al-Mīzān}, \emph{al-Kashshāf}, \emph{Majmaʿ al-Bayān}, and
\emph{al-Taḥrīr wa-l-Tanwīr}: \textbf{A} (anger-derived --- behavioural
completion of an inner-anger trajectory), \textbf{S} (structurally
derived --- greed/\emph{istikbār}/\emph{ʿulw} without explicit anger
antecedent), or \textbf{A+S} (joint). Coding was conducted by the
corresponding author with second-coder verification on a 20\% subsample
(κ = 0.79, substantial agreement; full coding at
\texttt{data/concordance/stage6\_causation\_coding.csv}). Results: A =
38 (26\%), S = 71 (49\%), A+S = 36 (25\%); by root, \emph{baghy}
22A/50S/24A+S, \emph{ṭughyān} 11A/19S/9A+S, \emph{ʿutuww} 5A/2S/3A+S.
Re-running the χ²(1) phenomenology-vs-outcomes split on the A-only
subset (Stages 1--5 = 167 vs.~anger-derived Stage-6 = 38) yields χ²(1) =
81.4 (\emph{p} \textless{} 0.001) --- the phenomenological core
dominates anger-derived outcomes \textasciitilde4.4×. The full-Stage-6
near-balance (χ²(1) = 1.55) is therefore explained: structurally-derived
attestations belong to an adjacent moral-anthropological field of
\emph{ẓulm-istikbār-ʿulw} that the corpus also pervasively documents,
and \emph{baghy}/\emph{ṭughyān}/\emph{ʿutuww} are not coextensive with
``anger outcomes'' in the strict sense.

\begin{figure}
\centering
\includegraphics[width=0.95\linewidth,height=\textheight,keepaspectratio]{fig4_cooccurrence.pdf}
\caption{Aya-level co-occurrence network of the fourteen core roots.
Each edge represents one or more verses in which two roots co-occur;
edge weight (on the line) is the count of shared ayat. The strongest
edge --- \emph{ḥzn--ḍyq}, weight 3 --- provides direct corpus-internal
validation of the Stage-2 (inner pressure \& contraction) grouping:
these two pre-anger lexemes share \emph{more} verses than any other pair
in the spectrum, indicating that the Qurʾān treats them as discursively
conjoint phenomena (``constriction of the chest in sadness,'' Q. 11:12;
Q. 6:33). \emph{Baghy} (Stage 6) anchors the right-hand cluster, with
edges reaching into the moral-evaluation lexicon (\emph{ẓlm},
\emph{jrm}, \emph{fsq}) --- a position formalised quantitatively in the
next figure. The transitive structure of the spectrum (no direct Stage-1
→ Stage-6 edges; mediation through Stages 2--5) is consistent with the
action-intensity gradient hypothesis, while also being consistent with
the bidirectional reading of Stage 6 developed in §4.8.5.}
\end{figure}

\begin{figure}
\centering
\includegraphics[width=0.95\linewidth,height=\textheight,keepaspectratio]{fig7_centrality.pdf}
\caption{Three centrality measures computed on the aya-level
co-occurrence graph: (a) weighted degree --- the total number of
co-occurrence ties; (b) betweenness --- how often a node lies on the
shortest path between others, indicating bridge-node strength; (c)
closeness --- the inverse mean distance to all other nodes. \emph{Baghy}
attains the \textbf{highest closeness centrality} (0.55) and shares the
highest betweenness centrality (0.15) with \emph{ghaḍab} and
\emph{asaf}, confirming its role as a principal \emph{pivot} between the
anger spectrum (Stages 1--5) and the wider moral-evaluation field of the
Qurʾān (315 attestations of \emph{ẓulm}, 66 of \emph{jurm}, 54 of
\emph{fisq}, 106 of \emph{ʿadāwah}). The closeness panel reveals a
complementary pattern: \emph{ḍyq} and \emph{ḥzn} (Stage 2) achieve high
closeness, indicating that the inner-pressure cluster is structurally
\emph{central} to the spectrum even though peripheral in raw frequency.
Together these measures support a layered reading of the network ---
Stage 2 supplies the affective core, Stage 6 supplies the structural
bridge to neighbouring moral-discursive fields, while Stages 3--5
mediate the energy transfer between them.}
\end{figure}

\subsubsection{4.9. Robustness checks: the umbrella moral
terms}\label{robustness-checks-the-umbrella-moral-terms}

For robustness, the umbrella roots were extracted and analysed (note
that \emph{karh}, formerly an umbrella term, is now incorporated into
the spectrum at Stage 1 and is reported in Table 2). Distributions of
the remaining five (Table 5) confirm that the spectrum's structure is
not artefactual.

\textbf{Table 5. Distribution of umbrella moral terms}

\begin{longtable}[]{@{}
  >{\centering\arraybackslash}p{(\linewidth - 10\tabcolsep) * \real{0.2727}}
  >{\centering\arraybackslash}p{(\linewidth - 10\tabcolsep) * \real{0.1364}}
  >{\centering\arraybackslash}p{(\linewidth - 10\tabcolsep) * \real{0.1364}}
  >{\centering\arraybackslash}p{(\linewidth - 10\tabcolsep) * \real{0.1364}}
  >{\centering\arraybackslash}p{(\linewidth - 10\tabcolsep) * \real{0.1364}}
  >{\raggedright\arraybackslash}p{(\linewidth - 10\tabcolsep) * \real{0.1818}}@{}}
\toprule\noalign{}
\begin{minipage}[b]{\linewidth}\centering
Root
\end{minipage} & \begin{minipage}[b]{\linewidth}\centering
Total
\end{minipage} & \begin{minipage}[b]{\linewidth}\centering
Meccan
\end{minipage} & \begin{minipage}[b]{\linewidth}\centering
Medinan
\end{minipage} & \begin{minipage}[b]{\linewidth}\centering
Binomial p
\end{minipage} & \begin{minipage}[b]{\linewidth}\raggedright
Interpretation
\end{minipage} \\
\midrule\noalign{}
\endhead
\bottomrule\noalign{}
\endlastfoot
šnʾ & 3 & 1 & 2 & --- & n too small for inference \\
jrm & 66 & 60 & 6 & \textless{} 0.0001 & Strongly Meccan;
anti-\emph{shirk} polemic \\
fsq & 54 & 20 & 34 & 0.0007 & Significantly Medinan;
community-discipline register \\
ẓlm & 315 & 211 & 104 & --- & Pervasive meta-category; cannot anchor a
single point \\
ʿdw & 106 & 47 & 59 & 0.008 & Tilts Medinan; the lexeme of mutual
antagonism \\
\end{longtable}

The opposing patterns of \emph{jrm} (strongly Meccan) and \emph{fsq}
(significantly Medinan) instantiate the same discursive-contextual
specialisation that we identified for \emph{sakhaṭ}. The umbrella
analysis thus confirms a \emph{general} phenomenon, of which the
\emph{sakhaṭ} finding is one specific instance.

\subsubsection{4.10. Comparative translation analysis (illustrative
sample)}\label{comparative-translation-analysis-illustrative-sample}

To illustrate how the spectrum's intensity differentiation fares in
translation, Table 6 compares five English translations on four
canonical verses. We frame this as an \emph{illustrative} case, not a
systematic translation-criticism evaluation: a properly powered
evaluation would require a stratified random sample of all 312
attestations (or, more practically, \textasciitilde30 verses sampled
across the six stages with inter-rater coding by trained Arabist
coders). The four-verse, five-translation sample here suffices to
register the \emph{pattern} of levelling but does not, on its own,
support the strong translation-recommendation set in §5.7; the
recommendations there are conditional on a follow-up evaluation that we
commend to future work.

\textbf{Table 6. Translation of four canonical verses across five
English translations} (illustrative sample; n = 4 verses × 5
translations = 20 cells)

\begin{longtable}[]{@{}
  >{\centering\arraybackslash}p{(\linewidth - 10\tabcolsep) * \real{0.2308}}
  >{\raggedright\arraybackslash}p{(\linewidth - 10\tabcolsep) * \real{0.1538}}
  >{\raggedright\arraybackslash}p{(\linewidth - 10\tabcolsep) * \real{0.1538}}
  >{\raggedright\arraybackslash}p{(\linewidth - 10\tabcolsep) * \real{0.1538}}
  >{\raggedright\arraybackslash}p{(\linewidth - 10\tabcolsep) * \real{0.1538}}
  >{\raggedright\arraybackslash}p{(\linewidth - 10\tabcolsep) * \real{0.1538}}@{}}
\toprule\noalign{}
\begin{minipage}[b]{\linewidth}\centering
Term \& verse
\end{minipage} & \begin{minipage}[b]{\linewidth}\raggedright
Yusuf Ali
\end{minipage} & \begin{minipage}[b]{\linewidth}\raggedright
Pickthall
\end{minipage} & \begin{minipage}[b]{\linewidth}\raggedright
Arberry
\end{minipage} & \begin{minipage}[b]{\linewidth}\raggedright
Saheeh Intl.
\end{minipage} & \begin{minipage}[b]{\linewidth}\raggedright
Hilali-Khan
\end{minipage} \\
\midrule\noalign{}
\endhead
\bottomrule\noalign{}
\endlastfoot
\emph{ḍīq} (Q. 15:97) & ``thy heart is distressed'' & ``thy bosom is at
times oppressed'' & ``thy breast is straitened'' & ``your chest is
constrained'' & ``your breast feels tight'' \\
\emph{ġayẓ} (Q. 3:134) & ``restrain anger'' & ``control their wrath'' &
``restrain their rage'' & ``suppress anger'' & ``repress anger'' \\
\emph{tamayyuz} (Q. 67:8) & ``almost burst with fury'' & ``almost burst
with rage'' & ``well-nigh bursts asunder with fury'' & ``almost bursts
with rage'' & ``almost bursts up with fury'' \\
\emph{ṭughyān} (Q. 96:6) & ``transgresses all bounds'' & ``rebelleth'' &
``waxes insolent'' & ``transgresses'' & ``transgresses (boundary)'' \\
\end{longtable}

Two illustrative observations emerge from this small sample. First, the
rendering of \emph{ġayẓ} in Q. 3:134 collapses the \emph{ghaḍab/ġayẓ}
distinction in all five translations---the English versions render
\emph{ġayẓ} with terms (\emph{anger}, \emph{wrath}, \emph{rage}) that
are also routinely used for \emph{ghaḍab}. The differentia---the
\emph{containment} dimension of \emph{ġayẓ}---is preserved only via the
contextual verb of restraint, not in the lexeme itself. Second, the
rendering of \emph{tamayyuz} in Q. 67:8 is uneven: Arberry's ``well-nigh
bursts asunder'' preserves the \emph{structural-disintegration}
component of the source; the others render it variationally as
\emph{burst with} + emotion, foregrounding the exit of contents but
losing the structural-collapse implication. These observations are
\emph{suggestive of} a systematic levelling pattern but are not on their
own evidentially decisive; we register them as motivating the systematic
translation evaluation we commend in §6.3.

\begin{center}\rule{0.5\linewidth}{0.5pt}\end{center}

\subsection{5. Discussion}\label{discussion}

\subsubsection{5.1. The integrated logic of
escalation}\label{the-integrated-logic-of-escalation}

The findings cohere into a single structural claim: the Qurʾānic lexicon
of the anger spectrum encodes a \textbf{logic of anger escalation} that
operates simultaneously along three dimensions---locus (inner → outer),
intensity (low → high), and scope (individual → structural). The six
stages instantiate this trifold transformation. The lexemes are not
synonyms; they are \emph{grades}, ordered by an intelligible cognitive
logic.

In compact form: \emph{Pre-anger displeasure → Inner pressure →
Evaluative aversion → Active anger → Compressed rage → Behavioural
outcomes (with caveats)}. The arrows index not mere temporal succession
but conditional intensification at each step where regulation may
succeed or fail. The Stage 6 step is, however, distinctive: the move
into behavioural outcomes is \emph{not} a deterministic escalation from
anger but is influenced by independent vectors (greed, \emph{istikbār},
structural arrogance), which is why we read the near-balance between
Stages 1--5 (167) and Stage 6 (145) as theoretically informative rather
than paradoxical.

\subsubsection{5.2. Convergence with contemporary emotion
psychology}\label{convergence-with-contemporary-emotion-psychology}

The Qurʾānic stages map onto Gross's (1998, 2014, 2015) process model of
emotion regulation, with the Qurʾānic spectrum partitioning more finely
on the affective side and adding an explicit behavioural-outcomes
register: Stage 1 (pre-anger \emph{uff/karh}) ↔ antecedent-focused
situation selection; Stage 2 (\emph{ḍīq/ḥuzn/asaf}) ↔ situation
modification; Stage 3 (\emph{naqm/sakhaṭ/maqt}) ↔ attentional deployment
/ cognitive change; Stage 4 (\emph{ghaḍab/ḥard}) ↔ response selection;
Stage 5 (\emph{ghayẓ}/\emph{kaẓm}; \emph{tamayyuz}) ↔ response
modulation and containment failure; Stage 6
(\emph{baghy/ṭughyān/ʿutuww}) --- where Gross's model becomes silent and
the Qurʾān adds the moral anthropology of structural-moral pathologies
that may be downstream from anger but, with bidirectional causation, may
also be upstream of it.

The Qurʾānic lexicon also realises an STAXI-style trichotomy
(Spielberger 1999): trait-like sustained states (\emph{ḥuzn},
\emph{ḍīq}), state-like situational responses (\emph{sakhaṭ},
\emph{maqt}, \emph{ghaḍab}), and deliberate restraint (\emph{kaẓm}) ---
three axes contemporary psychometrics has independently isolated.
Plutchik's (1980, 2001) annoyance → anger → rage three-level wheel maps
onto Stages 1--5 but stops at the regulatory boundary; Stage 6 is the
Qurʾān's substantive extension into the morally-anthropological.
Appraisal theory (Frijda 1986; Lazarus 1991; Roseman 1996; Scherer 2001,
2005) holds that emotions arise from cognitive evaluations rather than
raw stimuli; Q. 96:6--7 conditions \emph{ṭughyān} on the appraisal of
\emph{istighnāʾ} (perceived self-sufficiency) --- a textbook
appraisal-theoretic claim aligning with Roseman's ``agency-self /
motive-inconsistent / high-control'' pattern. The Qurʾānic lexicon thus
encodes --- as architectural design rather than explicit theoretical
statement --- insights contemporary emotion science has reached
independently.

\subsubsection{5.4. The phenomenology/outcomes
balance}\label{the-phenomenologyoutcomes-balance}

A naive ten-root reading once suggested 2.4× more attestations at the
behavioural pole than at Stages 1--3 combined, indicating discursive
over-emphasis on outcomes. The revised fourteen-root inventory (167 /
145; χ²(1) = 1.55, fail to reject) tells a different story: the expanded
phenomenological vocabulary (\emph{uff}, \emph{karh}, \emph{naqm},
\emph{maqt}, \emph{ḥard}) brings the lower stages into
corpus-distributional balance with Stage 6, exactly as one would expect
if (i) the phenomenological side is richer than the prior analysis
allowed, and (ii) Stage-6 lexemes are not exclusively anger-derived but
stand in bidirectional causal relation to anger. The methodological
lesson: corpus-distributional analyses of religious texts are maps of
\emph{discursive emphasis}, fundamentally entangled with
category-membership and category-causation decisions; the 10/4 → 14/6
shift changed both the empirical pattern and its interpretation.

\subsubsection{5.5. Discursive-contextual
specialisation}\label{discursive-contextual-specialisation}

The FDR-surviving Medinan exclusivity of \emph{sakhaṭ}, alongside the
contextual-only Meccan tilts of \emph{asaf}, \emph{naqm}, and
\emph{ʿutuww}, the hapax-Meccan \emph{ḥard}, the Medinan tilt of
\emph{karh}, and the strong umbrella patterns \emph{jrm} (Meccan) and
\emph{fsq} (Medinan), instantiate \textbf{discursive-contextual
specialisation}: members of multi-lexeme semantic fields preferentially
attaching to specific discursive registers. The pattern is not random
--- \emph{sakhaṭ} attaches to the \emph{Munāfiqūn} (a Medinan
phenomenon), \emph{asaf} to prior-prophet narratives
(Meccan-rhetorical), \emph{jrm} to anti-\emph{shirk} polemic (Meccan),
\emph{fsq} to intra-community discipline (Medinan). El-Awa (2006) on
textual relations and Sinai (2017, 2023) on Meccan/Medinan diachrony
supply the theoretical framework: specialisation is exactly the kind of
pattern one would expect from the Qurʾān's gradual response to evolving
sociolinguistic contexts. The phenomenon is, to our knowledge,
undocumented as a systematic corpus-distributional pattern in
Qurʾānic-linguistic scholarship; future work should test the
generalisation across other multi-lexeme fields, in dialogue with
QurAna/QurSim (Sharaf \& Atwell 2012) and the QuranMorph (2025) and
MASAQ extensions of the Dukes corpus.

\subsubsection{5.6. Q. 67:8 and the container
metaphor}\label{q.-678-and-the-container-metaphor}

Q. 67:8 (\emph{takādu tamayyazu min al-ġayẓ}) instantiates the ANGER IS
PRESSURIZED CONTAINER metaphor with unusual transparency: where English
idioms (\emph{explode}, \emph{blow one's top}, \emph{lose it}) report
the \emph{exit} of contained content, \emph{tamayyaza} (m-y-z: to part,
to sever) lexicalises \emph{the structural disintegration of the vessel
itself}. We frame this claim conservatively: assessing whether Q. 67:8
is the \emph{most} articulated instance in the historical record
requires a comparative survey of pre-modern Greek (Homeric
\emph{kholos}, \emph{menis}), Akkadian (\emph{ezzu}, \emph{libbāti
malû}), Sanskrit (\emph{krodha}, \emph{manyu}), Hebrew (\emph{ḥēmâ},
\emph{ʾaph}-idioms), and pre-Islamic \emph{Jāhilī} poetry --- a survey
beyond the scope of the present paper that we commend to specialists.
The Qurʾānic data are \emph{consistent with} CMT's universalist claim
that bodily-grounded image schemas (CONTAINER, FORCE) underwrite emotion
conceptualisation across cultures, while extending the empirical base to
seventh-century Arabic; they do not, on their own, establish the strong
cross-cultural universalism that earlier formulations of this paper
risked overstating. The two-volume \emph{Metaphors of ANGER across
Languages} (Kövecses, Benczes \& Szelid eds.~2025; De Gruyter),
conspicuously omitting Classical Arabic from its language coverage, is
the natural framework against which the present analysis should be
situated.

\begin{figure}
\centering
\includegraphics[width=1\linewidth,height=\textheight,keepaspectratio]{fig6_metaphor_diagram.pdf}
\caption{The ANGER-IS-A-PRESSURIZED-CONTAINER schema (Lakoff \& Kövecses
1987), instantiated by the Qurʾānic triplet \emph{ghaḍab → ghayẓ + kaẓm
→ tamayyuz}. \textbf{Panel 1} (Stage 4): \emph{ghaḍab} --- an
\emph{open} container; the affect is present but unconcentrated,
dissipating freely into language (the lexeme is the most nominally
productive in the spectrum). \textbf{Panel 2} (Stage 5, Q. 3:134):
\emph{ghayẓ} --- fluid sealed by \emph{kaẓm}; the verb \emph{kaẓama}
originally denotes the \emph{tying-off of a waterskin}, the
agent-internal regulation that maintains containment under rising
pressure. \textbf{Panel 3} (Stage 5 manifestation, Q. 67:8):
\emph{tamayyuz} --- the structural failure; the verb derives from the
root \emph{m-y-z} ``to separate, to part, to sever'', and figures not
the exit of pressurised content (as English \emph{explode}) but the
\emph{splitting apart of the container itself}. The Qurʾānic
instantiation is, on the present argument, more transparent than any
English exemplar in Lakoff and Kövecses's (1987) corpus, providing
strong cross-linguistic evidence for the cognitive universality of the
schema (§5.6).}
\end{figure}

\subsubsection{5.7. Implications for Qurʾānic translation (conditional
on follow-up
evaluation)}\label{implications-for-qurux2beux101nic-translation-conditional-on-follow-up-evaluation}

Translation criticism of the Qurʾān has long observed that translations
tend to level the spectrum---that \emph{ġaḍab} and \emph{ġayẓ} alike
become \emph{anger}; that \emph{baghy}, \emph{ṭughyān}, and
\emph{ʿutuww} alike become \emph{transgression} (Table 6). Our
four-verse, five-translation illustrative sample (§4.10) is consistent
with this pattern but is too small (\emph{n} = 4 verses × 5 translations
= 20 cells) to support strong prescriptive claims; recent computational
work --- notably Gaanoun \& Alsuhaibani (2025; \emph{Nature HSSC}) on
sentiment preservation across seven English translations and Alsuhaibani
et al.'s (2025) ELQV emotion-labelled dataset of 2,100 Qurʾānic verses
--- has begun to operationalise the levelling concern at scale on coarse
Ekman categories, though without a graded reference model of the lexical
spectrum against which the levelling could be benchmarked. We frame the
present discussion as \textbf{conditional implications}: contingent on a
follow-up systematic evaluation (a stratified random sample of
\textasciitilde30 verses across the six stages, with multi-rater coding
against the spectrum reference model), the analysis would \emph{if
confirmed} support the following candidate equivalents for critical and
study-edition translations:

\begin{quote}
\emph{uff}: vocal sigh of irritation · \emph{karh} (\emph{karāha}):
inner dislike · \emph{karh} (\emph{ikrāh}): coercion (a structural
extension, not an emotion) · \emph{ḍīq}: inner constriction · \emph{ḍāqa
dharʿan}: arm-span tightening / inhibited aggression · \emph{ḥuzn}:
persistent grief (transversal valence) · \emph{asaf}: regretful dismay
(with anger) · \emph{naqm}: vengeful disapproval · \emph{sakhaṭ}: moral
disapprobation · \emph{maqt}: intense moral aversion · \emph{ghaḍab}:
anger · \emph{ḥard}: anger combined with deliberate denial ·
\emph{ghayẓ}: contained / boiling rage · \emph{tamayyuz min al-ġayẓ}:
bursting from rage · \emph{baghy}: aggressive transgression ·
\emph{ṭughyān}: rebellious overflow · \emph{ʿutuww}: obstinate defiance.
\end{quote}

A second conditional implication, which we register as a hypothesis
rather than a recommendation, is that pedagogical translations may
benefit from marking each lexeme's position on the continuum (by inline
gloss, footnote, or chromatic coding) --- but the empirical case for
this rests on the comparative comprehension study we commend to future
work. The 30-verse stratified evaluation, paired with an inter-rater
coded comparison across the present spectrum reference model and the
ELQV labels, is a natural Phase II for the present dataset.

\subsubsection{5.8. Implications for Islamic
psychology}\label{implications-for-islamic-psychology}

The continuum offers a culturally-grounded framework mapping onto
Gross's model without importation of secular vocabulary, with three
conceptual implications: \textbf{(i) Prevention} --- early-stage
recognition through indigenous Stages 1--2 lexicon (\emph{uff},
\emph{karāha}, \emph{ḍīq}, \emph{ḥuzn}) supports CBT-style interception
of affective dysregulation; \textbf{(ii) Crisis intervention} --- the
\emph{kaẓm} construct integrates with contemporary
suppression-and-reappraisal techniques; \textbf{(iii) Rehabilitation}
--- the post-failure Stage 6 is framed as recoverable through
\emph{tawba}, providing motivational reinforcement for anger-disorder
and post-violent-trauma rehabilitation. These implications are
conceptual; their translation into validated clinical interventions
awaits empirical studies we recommend as a priority for future work.

\begin{center}\rule{0.5\linewidth}{0.5pt}\end{center}

\subsection{6. Conclusion}\label{conclusion}

\subsubsection{6.1. Substantive
contributions}\label{substantive-contributions}

This paper has advanced three contributions to the scholarship.

\textbf{Methodologically}, we have demonstrated that the integration of
three previously separate frameworks---Izutsu's semantic-field analysis,
Kennedy--Sassoon scalar semantics, and Lakoff--Kövecses
conceptual-metaphor theory---coupled with a reproducible computational
pipeline over the Quranic Arabic Corpus, produces a research apparatus
stronger than the sum of its parts. To our knowledge, this trifold
integration with corpus-empirical validation has not previously been
deployed in Qurʾānic studies as a unified analytic framework, although
Qāʾimīniyā (1390/2011, 1393/2014) and Pakatchi have integrated Izutsu
and conceptual-metaphor theory, and Persian-language scalar-semantic
work on Qurʾānic vocabulary may exist in venues we have not been able to
canvass. The research apparatus is openly available and applicable to
other Qurʾānic semantic fields.

\textbf{Descriptively}, we have produced an exhaustive, reproducible
concordance of 312 attestations across the fourteen core spectrum roots,
with corresponding distributional and network-analytic outputs. The CSV
outputs --- including FDR-adjusted binomial tests, marginal-preserving
permutation null \emph{p}-values, and bootstrap 95\% CIs on every
centrality estimate --- provide a foundation on which further
philological and computational work can build.

\textbf{Theoretically}, we have formulated and empirically defended the
\textbf{Qurʾānic Action-Intensity Continuum Hypothesis}: the claim that
fourteen core lexemes of the Qurʾānic anger-spectrum field operate as a
single graded continuum structured by \emph{intensity of action} and
partitioned into six phenomenologically meaningful stages. The
hypothesis is supported by qualitative evidence from the four
authoritative lexica and from a now-broadened tafsir consultation
(§3.5), and by quantitative evidence from frequency, distribution,
morphology, and network structure. We have, in addition, identified
distributional patterns that, to our knowledge, have not been previously
reported as a coordinated set, while reporting their post-FDR evidential
strength transparently: the exclusively-Medinan \emph{sakhaṭ} survives
Holm-Bonferroni correction at α = 0.05 \emph{borderline} (raw \emph{p} =
0.005, adjusted \emph{p} = 0.05) and is robust across baselines
0.70--0.78; the exclusively-Meccan \emph{asaf}, the Meccan-tilted
\emph{naqm}, the hapax-Meccan \emph{ḥard}, and the Medinan-tilted
\emph{karh} are reported as \emph{contextual observations} warranting
follow-up rather than as confirmed findings. The general phenomenon of
\textbf{discursive-contextual lexical specialisation} of which we
hypothesise these are instances is proposed as a research target rather
than as a settled result. We have entered explicit caveats on the
bidirectional causation of the Stage-6 lexemes (§4.8.5), grounded in a κ
= 0.79-validated tafsir-coding audit, and argued that Q. 67:8 furnishes
an \emph{unusually transparent} linguistic exemplar of ``container
collapse''---a comparative claim whose strong universalist form is
reserved pending the cross-cultural philological survey we commend in
§5.6 and §6.3.

\subsubsection{6.2. Answers to the research
questions}\label{answers-to-the-research-questions}

\textbf{RQ1.} The semantic field is a fourteen-node network organised
around the bridge node \emph{baghy} and the stage-internal hubs
\emph{ḥuzn} and \emph{ghaḍab}, with several specialised peripheral nodes
(\emph{sakhaṭ}, \emph{asaf}, \emph{ḥard}, \emph{ʿutuww}) exhibiting
marked distributional asymmetries.

\textbf{RQ2.} The fourteen nodes order on a six-stage
continuum---pre-anger displeasure (Stage 1), inner pressure \&
contraction (Stage 2), evaluative aversion (Stage 3), active anger
(Stage 4), compressed/explosive rage (Stage 5), behavioural outcomes
with caveats (Stage 6)---organised by the composite dimension of
\emph{intensity of action} (locus, intensity, scope of consequence).

\textbf{RQ3.} The empirical distribution is \emph{consistent with}
(though not by itself confirmatory of) the continuum: the asymptotic
chi-squared rejects the null of six-stage uniform distribution (χ²(5) =
227.15; large effect size at Cohen's \emph{w} = 0.85), but a
marginal-preserving permutation null on the same statistic yields
\emph{p} = 0.24, indicating that the asymptotic \emph{p}-value
over-states the evidence on this sparse, marginally-imbalanced design;
the substantive case for the six-stage architecture rests on the
qualitative lexicographic, exegetical, and metaphorical evidence rather
than on the omnibus \emph{p}-value alone. The comparison between Stages
1--5 and Stage 6 \emph{fails to reject} equal split (χ²(1) = 1.55,
Cohen's \emph{w} = 0.07, trivial effect); we read this as
\emph{consistent with} --- but not as positive evidence for --- the
bidirectional-causation reading of Stage 6, with the substantive support
coming from the κ = 0.79-validated tafsir-coding audit (§4.8.5). The
network exhibits a bridging-role profile at \emph{baghy} with wide
bootstrap confidence intervals reported transparently. The morphological
profile distinguishes the verbal Stages 1--2 from the more nominal Stage
6.

\textbf{RQ4.} The continuum is structurally homologous with Gross's
process model, Plutchik's wheel, Spielberger's STAXI, and
Lazarus--Scherer appraisal theory in its psychological component, but
extends beyond all of them by integrating the post-failure
structural-moral consequences of unmanaged emotion (Stage 6) --- while
qualifying that integration with the bidirectional-causation caveat of
§4.8.5.

\textbf{RQ5.} \emph{Conditional} practical implications include
candidate distinct-equivalent translations and position-marking
conventions (subject to a 30-verse stratified follow-up evaluation,
§5.7) and three Islamic-psychology applications (early-stage prevention,
\emph{kaẓm}-grounded crisis intervention, \emph{tawba}-grounded
rehabilitation). All practical implications are framed as hypotheses for
empirical follow-up, not as settled prescriptions.

\subsubsection{6.3. Limitations and future
work}\label{limitations-and-future-work}

Five limitations bound the study: (i) surface-lexical scope ---
discourse-analytic close-reading of the spectrum's narrative deployment
across suras is left for future work; (ii) the broadened nine-tafsir
engagement is still selective, with deeper Sufi (al-Qushayrī, Ibn
ʿArabī's school) and Persian-language commentary needing fuller
integration; (iii) the cognitive-linguistic apparatus is
Anglophone-dominated --- engagement with classical Arabic \emph{bayān}
theory (al-Jurjānī's \emph{Dalāʾil al-Iʿjāz}) would supply an indigenous
substrate; (iv) the \emph{ḥadīth} corpus has not been integrated; (v)
the §5.8 Islamic-psychology implications await operationalisation in
pilot studies of \emph{kaẓm}-grounded emotion-regulation training. A
natural follow-up is the AraBERT/CAMeLBERT contextual-embedding approach
to the same fourteen-root inventory (testing stage-recoverability from
distributional structure), deferred here to maintain the
philological-statistical centre of gravity.

\subsubsection{6.4. Closing remark}\label{closing-remark}

The Qurʾānic lexicon encodes a sophisticated theory of the anger
spectrum --- not as explicit theoretical formulation but as the
architectural design of its semantic surface. Recovering that theory
requires methodological tools that bring together classical philological
attentiveness, contemporary linguistic-theoretical precision, and
corpus-computational verifiability. We hope the present study, alongside
Izutsu, Qāʾimīniyā, Pakatchi, Narimani et al., and Qarizadeh et al.,
advances this triple integration. The classical and the computational
are complementary instruments whose joint deployment opens horizons
neither could disclose alone.

\begin{center}\rule{0.5\linewidth}{0.5pt}\end{center}

\subsection{Acknowledgements}\label{acknowledgements}

The authors thank Kais Dukes (in memoriam) and the team behind the
Quranic Arabic Corpus, and the Tanzil Project, for releasing the
morphological and textual resources without which this kind of empirical
Qurʾānic analysis would not be feasible. We thank Urmia University for
institutional support, and the anonymous peer reviewer whose detailed
comments materially improved the methodological framing, the statistical
interpretation (especially regarding the χ²(1) result and the
\emph{sakhaṭ} baseline-sensitivity question), and the reproducibility of
the analysis pipeline.

\subsection{Funding}\label{funding}

This research received no specific grant from any funding agency in the
public, commercial, or not-for-profit sectors.

\subsection{Competing interests}\label{competing-interests}

The authors declare no competing financial or non-financial interests.

\subsection{Ethics statement}\label{ethics-statement}

This study uses only publicly-released text-corpus data (the Quranic
Arabic Corpus, GPL-licensed; the Tanzil Quran text, CC-BY-ND 3.0). No
human or animal subjects were involved at any stage. As a corpus-only
philological-computational analysis, the work is exempt from
human-subjects review under the standards of Urmia University and of
comparable institutions.

\subsection{Author contributions}\label{author-contributions}

K.J. (corresponding author) conceived the project, developed the
theoretical framework, conducted the lexicographic and tafsir-tradition
analysis, drafted the manuscript, and is responsible for the final form
of the argument. A.J. contributed to the theoretical framework, the
design of the computational pipeline, and editorial revision. Both
authors approved the final manuscript and are jointly accountable for
its content.

\subsection{Author affiliations and
ORCID}\label{author-affiliations-and-orcid}

\textbf{Karim Jabbary} (corresponding author). Department of Arabic
Language and Literature, Faculty of Literature and Humanities, Urmia
University, Urmia, Iran. \emph{email:} k.jabbary@urmia.ac.ir.

\textbf{Ali Jabbary}. Faculty of Computer Engineering, Urmia University,
Urmia, Iran. \emph{email:} a.jabbary@urmia.ac.ir.

(ORCID identifiers will be provided at submission.)

\subsection{Data and Code
Availability}\label{data-and-code-availability}

All code, data, concordance CSV outputs, and figure-generation scripts
are released open-source at:

\begin{quote}
\href{https://github.com/ali-kin4/quranic-emotion-spectrum}{github.com/ali-kin4/quranic-emotion-spectrum}
\end{quote}

The repository contains: (i) the Python pipeline under MIT License; (ii)
the Quranic Arabic Corpus (Dukes 2011) under its original GPL; (iii) the
complete CSV concordance for the fourteen core spectrum roots (312
attestations), distributional summaries, statistical tests (χ²(5) =
227.15; χ²(1) = 1.55), network-centrality outputs, and
umbrella-cooccurrence tables; (iv) PDF and PNG versions of the seven
figures published with this paper; (v) full text of the Persian and
English manuscripts; (vi) the SHA-256 hash of the QAC v0.4 source file
(\texttt{data/quran/qac\_morphology.sha256}) and the manual-validation
audit (\texttt{data/concordance/validation\_audit.md}). The complete
analysis can be regenerated end-to-end by:

\begin{Shaded}
\begin{Highlighting}[]
\ExtensionTok{pip}\NormalTok{ install }\AttributeTok{{-}r}\NormalTok{ analysis/requirements.txt}
\ExtensionTok{python}\NormalTok{ analysis/extract\_concordance.py}
\ExtensionTok{python}\NormalTok{ analysis/advanced\_metrics.py}
\ExtensionTok{python}\NormalTok{ analysis/visualize\_spectrum.py}
\ExtensionTok{python}\NormalTok{ analysis/visualize\_advanced.py}
\end{Highlighting}
\end{Shaded}

\begin{center}\rule{0.5\linewidth}{0.5pt}\end{center}

\subsection{References}\label{references}

Akra, D., Hammouda, T., \& Jarrar, M. (2025). \emph{QuranMorph:
Morphologically Annotated Quranic Corpus} {[}Preprint{]}.
arXiv:2506.18148. https://arxiv.org/abs/2506.18148

Albayrak, I. (2020). Semantics of the Qurʾānic Weltanschauung: A
critical analysis of Toshihiko Izutsu's works. \emph{American Journal of
Islam and Society}, \emph{37}(1--2).
https://ajis.org/index.php/ajiss/article/view/387

Almuhanna, M., et al.~(2025). Sentiment preservation in Quran
translation with artificial intelligence approach: Study in reputable
English translation of the Quran. \emph{Humanities and Social Sciences
Communications}, \emph{12}(1), Article 04181.
https://doi.org/10.1057/s41599-024-04181-0

Averill, J. R. (1982). \emph{Anger and Aggression: An Essay on Emotion}.
New York: Springer-Verlag.

Berkowitz, L. (1993). \emph{Aggression: Its Causes, Consequences, and
Control}. New York: McGraw-Hill.

Derki, M. (2022). Conceptualization of anger in Modern Standard Arabic
and English: A comparative study. \emph{Professional Discourse \&
Communication}. https://www.pdc-journal.com/jour/article/view/168

Dukes, K. (2011). \emph{Quranic Arabic Corpus (Morphology, Version
0.4)}. University of Leeds. https://corpus.quran.com

El-Awa, S. M. S. (2006). \emph{Textual Relations in the Qur'an:
Relevance, Coherence and Structure}. London and New York: Routledge.

Frijda, N. H. (1986). \emph{The Emotions}. Cambridge and Paris:
Cambridge University Press \& Éditions de la Maison des Sciences de
l'Homme.

Gross, J. J. (1998). The emerging field of emotion regulation: An
integrative review. \emph{Review of General Psychology}, \emph{2}(3),
271--299.

Gross, J. J. (Ed.). (2014). \emph{Handbook of Emotion Regulation} (2nd
ed.). New York: Guilford Press.

Gross, J. J. (2015). Emotion regulation: Current status and future
prospects. \emph{Psychological Inquiry}, \emph{26}(1), 1--26.

Ibn Fāris, A. (1979). \emph{Muʿjam Maqāyīs al-Lughah} (ʿA. M. Hārūn,
Ed.). Beirut: Dār al-Fikr.

Ibn Manẓūr, M. (1993). \emph{Lisān al-ʿArab} (3rd ed.). Beirut: Dār
Ṣādir.

Ibn ʿĀshūr, M. T. (1984). \emph{al-Taḥrīr wa-l-Tanwīr}. Tunis: al-Dār
al-Tūnisiyyah li-l-Nashr.

Izutsu, T. (1959). \emph{The Structure of the Ethical Terms in the
Koran: A Study in Semantics}. Tokyo: Keio Institute of Philological
Studies.

Izutsu, T. (1964). \emph{God and Man in the Koran: Semantics of the
Koranic Weltanschauung}. Tokyo: Keio Institute of Cultural and
Linguistic Studies.

Izutsu, T. (1966/2002). \emph{Ethico-Religious Concepts in the Qurʾān}.
Montreal: McGill-Queen's University Press.

Kassinove, H., \& Tafrate, R. C. (2002). \emph{Anger Management: The
Complete Treatment Guidebook for Practitioners}. Atascadero, CA: Impact
Publishers.

Kennedy, C. (1999). \emph{Projecting the Adjective: The Syntax and
Semantics of Gradability and Comparison}. New York: Garland Publishing.
(Published version of the 1997 UCSC PhD dissertation.)

Kennedy, C. (2007). Vagueness and grammar: The semantics of relative and
absolute gradable adjectives. \emph{Linguistics and Philosophy},
\emph{30}(1), 1--45.

Kövecses, Z. (1986). \emph{Metaphors of Anger, Pride, and Love: A
Lexical Approach to the Structure of Concepts}. Amsterdam: John
Benjamins.

Kövecses, Z. (2000). \emph{Metaphor and Emotion: Language, Culture, and
Body in Human Feeling}. Cambridge: Cambridge University Press.

Kövecses, Z. (2010). \emph{Metaphor: A Practical Introduction} (2nd
ed.). Oxford: Oxford University Press.

Lakoff, G., \& Johnson, M. (1980). \emph{Metaphors We Live By}. Chicago:
University of Chicago Press.

Lakoff, G., \& Kövecses, Z. (1987). The cognitive model of anger
inherent in American English. In D. Holland \& N. Quinn (Eds.),
\emph{Cultural Models in Language and Thought} (pp.~195--221).
Cambridge: Cambridge University Press.

Lazarus, R. S. (1991). \emph{Emotion and Adaptation}. New York: Oxford
University Press.

Maalej, Z. (2004). Figurative language in anger expressions in Tunisian
Arabic: An extended view of embodiment. \emph{Metaphor and Symbol},
\emph{19}(1), 51--75.

Maalej, Z. A., \& Yu, N. (Eds.). (2011). \emph{Embodiment via Body
Parts: Studies from Various Languages and Cultures}. Amsterdam and
Philadelphia: John Benjamins. https://doi.org/10.1075/hcp.31

al-Muṣṭafawī, Ḥ. (1416 AH / 1995). \emph{al-Taḥqīq fī Kalimāt al-Qurʾān
al-Karīm} (14 vols.). Tehran: Wizārat al-Thaqāfah wa-l-Irshād al-Islāmī.

Narimani, Z., Iqbali, M., \& Chari, M. (1400 SH / 2021). {[}Semantic
re-analysis of the words \emph{ghayẓ}, \emph{ghaḍab}, and \emph{sakhaṭ}
in the Holy Qurʾān with an exegetical approach{]}. \emph{Pizhūhish-hā-yi
Qurʾān va Ḥadīth} {[}Qurʾanic and Hadith Research{]}. University of
Tehran. {[}in Persian; original Persian title: «بازشناسی معنایی واژگان
«غیظ»، «غضب» و «سَخَط» در قرآن کریم با رویکرد تفسیری»{]}

Pakatchi, A., \& Afrashi, A. (1399 SH / 2020). \emph{Rūykardhā-yi
Maʿnā-shinākhtī dar Muṭālaʿāt-i Qurʾānī} {[}Semantic Approaches in
Qurʾānic Studies{]}. Tehran: Pizhūhishgāh-i ʿUlūm-i Insānī va
Muṭālaʿāt-i Farhangī. {[}in Persian{]}

Plutchik, R. (1980). \emph{Emotion: A Psychoevolutionary Synthesis}. New
York: Harper \& Row.

Plutchik, R. (2001). The nature of emotions. \emph{American Scientist},
\emph{89}(4), 344--350.

Qarizadeh, M. ʿO., Salmani Marvast, M. ʿA., Meymandi, V., \& Farsi, B.
(1402 SH / 2024). {[}Polysemy of the word \emph{ṭughyān} in the Holy
Qurʾān with a linguistic approach{]}. \emph{Pizhūhish-hā-yi
Zabānshinākhtī-yi Qurʾān} {[}Linguistic Research in the Holy Qurʾān{]}.
University of Isfahan. {[}in Persian; original Persian title: «چندمعنایی
واژۀ «طغیان» در قرآن کریم با رویکردی بر زبان‌شناختی»{]}

Qāʾimīniyā, ʿA. (1390 SH / 2011). \emph{Maʿnāshināsī-yi Shinākhtī-yi
Qurʾān} {[}Cognitive Semantics of the Qurʾān{]}. Tehran: Pizhūhishgāh-i
Farhang va Andīshah-i Islāmī. {[}in Persian; original Persian title:
«معناشناسی شناختی قرآن»{]}

Qāʾimīniyā, ʿA. (1393 SH / 2014). \emph{Istiʿārah-hā-yi Mafhūmī wa
Faḍāhā-yi Qurʾānī} {[}Conceptual Metaphors and Qurʾānic Spaces{]}.
Tehran: Pizhūhishgāh-i Farhang va Andīshah-i Islāmī. {[}in Persian;
original Persian title: «استعاره‌های مفهومی و فضاهای قرآنی»{]}

Rabab'ah, K., \& Al-Saidat, E. (2022). The conceptualization of anger
through metaphors, metonymies and metaphtonymies in Jordanian Arabic and
English: A contrastive study. \emph{Cognitive Semantics}, \emph{8}(3),
409--433. https://doi.org/10.1163/23526416-bja10044

Roseman, I. J. (1996). Appraisal determinants of emotions: Constructing
a more accurate and comprehensive theory. \emph{Cognition and Emotion},
\emph{10}(3), 241--277. https://doi.org/10.1080/026999396380240

al-Rāghib al-Iṣfahānī, Ḥ. (1412 AH / 1992). \emph{al-Mufradāt fī Gharīb
al-Qurʾān} (Ṣ. ʿA. Dāwūdī, Ed.). Damascus and Beirut: Dār al-Qalam \&
al-Dār al-Shāmiyyah.

Sassoon, G. W. (2010). The degree functions of negative adjectives.
\emph{Natural Language Semantics}, \emph{18}(2), 141--181.

Sawalha, M., Al-Shargi, F., Yagi, S., AlShdaifat, A. T., Hammo, B.,
Belajeed, M., \& Al-Ogaili, L. R. (2024). Morphologically-analyzed and
syntactically-annotated Quran dataset (MASAQ). \emph{Data in Brief},
\emph{58}, Article 111211. https://doi.org/10.1016/j.dib.2024.111211

Scherer, K. R. (2001). Appraisal considered as a process of multilevel
sequential checking. In K. R. Scherer, A. Schorr, \& T. Johnstone
(Eds.), \emph{Appraisal Processes in Emotion: Theory, Methods, Research}
(pp.~92--120). New York: Oxford University Press.

Scherer, K. R. (2005). What are emotions? And how can they be measured?
\emph{Social Science Information}, \emph{44}(4), 695--729.

Sharaf, A. M., \& Atwell, E. (2012a). QurAna: Corpus of the Quran
annotated with pronominal anaphora. In \emph{Proceedings of the Eighth
International Conference on Language Resources and Evaluation (LREC'12)}
(pp.~130--137). Istanbul: ELRA. https://aclanthology.org/L12-1011/

Sharaf, A. M., \& Atwell, E. (2012b). QurSim: A corpus for evaluation of
relatedness in short texts. In \emph{Proceedings of the Eighth
International Conference on Language Resources and Evaluation (LREC'12)}
(pp.~2295--2302). Istanbul: ELRA. https://aclanthology.org/L12-1051/

Sharifian, F. (2011). \emph{Cultural Conceptualisations and Language:
Theoretical Framework and Applications}. Amsterdam: John Benjamins.

Sinai, N. (2017). \emph{The Qur'an: A Historical-Critical Introduction}.
Edinburgh: Edinburgh University Press.

Sinai, N. (2023). \emph{Key Terms of the Qur'an: A Critical Dictionary}.
Princeton, NJ: Princeton University Press.

Soriano, C. (2003). Some anger metaphors in Spanish and English: A
contrastive review. \emph{International Journal of English Studies
(IJES)}, \emph{3}(2), 107--122.

Soriano, C. (2005). \emph{The Conceptualization of Anger in English and
Spanish: A Cognitive Approach} {[}Doctoral dissertation, Universidad de
Murcia{]}. https://digitum.um.es/digitum/handle/10201/33082

Soriano, C. (2013). Anger metaphors across languages: A cognitive
linguistic perspective. In R. R. Heredia \& A. B. Cieślicka (Eds.),
\emph{Bilingual Figurative Language Processing} (pp.~410--440).
Cambridge: Cambridge University Press.

Spielberger, C. D. (1999). \emph{Professional Manual for the State-Trait
Anger Expression Inventory--2 (STAXI-2)}. Odessa, FL: Psychological
Assessment Resources.

al-Ṭabarī, M. ibn J. (n.d. / 2000 ed.). \emph{Jāmiʿ al-Bayān ʿan Taʾwīl
Āy al-Qurʾān} (A. M. Shākir, Ed., 24 vols.). Beirut: Muʾassasat
al-Risālah.

al-Ṭabāṭabāʾī, M. Ḥ. (1390 AH / 1970). \emph{al-Mīzān fī Tafsīr
al-Qurʾān} (20 vols.). Beirut: Muʾassasat al-Aʿlamī li-l-Maṭbūʿāt.

al-Ṭabarsī, F. al-Ḥ. (n.d.). \emph{Majmaʿ al-Bayān fī Tafsīr al-Qurʾān}.
Beirut: Dār al-Maʿrifah.

\emph{Tanzil Project}. (n.d.). \emph{The Tanzil Quran Text} (Uthmani
edition). https://tanzil.net.

al-Qurṭubī, M. ibn A. (1964). \emph{al-Jāmiʿ li-Aḥkām al-Qurʾān} (20
vols.). Cairo: Dār al-Kutub al-Miṣriyyah.

al-Qushayrī, ʿA. al-K. (2000). \emph{Laṭāʾif al-Ishārāt} (I. Basūnī,
Ed., 6 vols.). Cairo: al-Hayʾah al-Miṣriyyah al-ʿĀmmah li-l-Kitāb.

Quṭb, S. (1980). \emph{Fī Ẓilāl al-Qurʾān} (6 vols.). Cairo and Beirut:
Dār al-Shurūq.

Riḍā, M. R. (1947). \emph{Tafsīr al-Qurʾān al-Ḥakīm (al-Manār)} (12
vols.). Cairo: Dār al-Manār.

Ibn Kathīr, I. (1419 AH / 1999). \emph{Tafsīr al-Qurʾān al-ʿAẓīm} (S.
al-Salāmah, Ed., 8 vols.). Riyadh: Dār Ṭaybah.

Wierzbicka, A. (1999). \emph{Emotions Across Languages and Cultures:
Diversity and Universals}. Cambridge: Cambridge University Press.

Yu, N. (1995). Metaphorical expressions of anger and happiness in
English and Chinese. \emph{Metaphor and Symbolic Activity},
\emph{10}(2), 59--92. https://doi.org/10.1207/s15327868ms1002\_1

al-Zamakhsharī, M. (1407 AH). \emph{al-Kashshāf ʿan Ḥaqāʾiq Ghawāmiḍ
al-Tanzīl}. Beirut: Dār al-Kitāb al-ʿArabī.
